\pdfminorversion=7
\documentclass[a4paper,final]{article}
\usepackage[14pt]{extsizes}
\usepackage[T2A]{fontenc}
\usepackage[utf8]{inputenc}
\usepackage[russian]{babel}

\addto\captionsrussian{\renewcommand{\figurename}{Рис.}}

% --- Math ---
\usepackage{amsmath,amssymb,amsfonts}
\usepackage{nccmath}

% --- Page geometry ---
\usepackage[left=20mm, top=20mm, right=20mm, bottom=20mm, footskip=10mm]{geometry}

% --- Text/layout helpers ---
\usepackage{ragged2e}
\usepackage{setspace}
\usepackage{indentfirst}

% --- Graphics ---
\usepackage{graphicx}
\DeclareGraphicsExtensions{.jpg,.png,.pdf}
\usepackage{float}
\usepackage{wrapfig}
\usepackage{tikz}
\usepackage{pgfplots}
\pgfplotsset{compat=1.18}

% --- Captions ---
\usepackage[font=small,singlelinecheck=false,justification=centering,labelsep=period]{caption}

% --- Lists ---
\usepackage{enumitem}

% --- Hyperlinks ---
\usepackage{hyperref}
\hypersetup{colorlinks=true, linkcolor=blue, urlcolor=blue}

% --- Code listings ---
\usepackage{xcolor}
\usepackage{listings}
\usepackage{pdfpages}
\definecolor{codegreen}{rgb}{0,0.6,0}
\definecolor{codegray}{rgb}{0.5,0.5,0.5}
\definecolor{codepurple}{rgb}{0.58,0,0.82}
\definecolor{backcolour}{rgb}{0.95,0.95,0.92}

\lstdefinestyle{mystyle}{
	backgroundcolor=\color{backcolour},
	commentstyle=\color{codegreen},
	keywordstyle=\color{magenta},
	numberstyle=\tiny\color{codegray},
	stringstyle=\color{codepurple},
	basicstyle=\ttfamily\footnotesize,
	breaklines=true,
	breakatwhitespace=false,
	captionpos=b,
	keepspaces=true,
	numbers=left,
	numbersep=5pt,
	showspaces=false,
	showstringspaces=false,
	showtabs=false,
	tabsize=4,
	columns=fullflexible,
	extendedchars=true,
	inputencoding=utf8,
	literate=
	{•}{{-}}1
	{─}{{-}}1
	{–}{{-}}1
	{—}{{-}}1
	{→}{{$\to$}}1
	{≠}{{$\neq$}}1
	{³}{{$^3$}}1
	{ü}{{\"u}}1
}

\lstset{style=mystyle}
\renewcommand{\lstlistingname}{Листинг}
\newcommand{\coderef}[2]{\hyperref[#2]{\texttt{\detokenize{#1}}}}
\DeclareUnicodeCharacter{202F}{\,}
\sloppy

% --- Other document tweaks ---
\renewcommand{\arraystretch}{1.4}
\parindent=24pt
\parskip=0pt
\tolerance=2000
\flushbottom

\addto\captionsrussian{\def\refname{Список использованной литературы}}

\begin{document}
	\begin{center}
		\normalsize\bf{{МИНИСТЕРСТВО НАУКИ И ВЫСШЕГО ОБРАЗОВАНИЯ РОССИЙСКОЙ ФЕДЕРАЦИИ}

			\hfill \break
			федеральное государственное автономное образовательное учреждение высшего образования\\ «Санкт-Петербургский политехнический университет\\ Петра Великого»\\[10pt]}
		\hfill \break
		\normalsize{Институт компьютерных наук и кибербезопасности}\\[10pt]

		\normalsize{Высшая школа технологий и искусственного интеллекта}\\[10pt]

		\normalsize{Направление: 02.03.01 Математика и компьютерные науки}\\
		\vspace{50pt}

		\vspace{1pt}
		\begin{center}
Отчет по курсовой работе\\
по дисциплине «Теория графов»\\
Реализация алгоритмов теории графов\\
для лабораторных работ №1--5

		\end{center}

		\hfill \break
		\hfill \break
	\end{center}
	\small{ \bf
		\begin{tabular}{lrrl}
			\!\!\!Студент, & \hspace{0cm} & & \\
			\!\!\!группа 5130201/40001 & \hspace{1.5cm} & \underline{\hspace{3cm}} & \hspace{-0.3cm}Ови Мд Шамин Ясир  \\\\
			\!\!\!Преподаватель, & \hspace{1cm} & \\
			\!\!\! & \hspace{0.1cm} & \underline{\hspace{3cm}} &  Востров Алексей Владимирович\\\\
			&&\hspace{3.5cm}
		\end{tabular}
		\begin{flushright}
			<<\underline{\hspace{1cm}}>>\underline{\hspace{2.5cm}} 2026г.
		\end{flushright}
	}

	\hfill \break
	\hfill \break
	\begin{center} \small\bf{Санкт-Петербург, 2026} \end{center}
	\thispagestyle{empty}
	% END OF TITLE PAGE

\newpage
\renewcommand{\contentsname}{Содержание}
\tableofcontents


\newpage
\section*{Введение}
\addcontentsline{toc}{section}{Введение}

\indent Отчёт состоит из описания 5 лабораторных работ. Каждая работа включает в себя реализацию различных алгоритмов на случайно созданных связных ациклических графах. Генерация графов осуществляется с использованием \textit{нормального распределения} (по справочнику Вадзинского~\cite{vadzinsky}).

\medskip
\noindent\textbf{ЛАБОРАТОРНАЯ РАБОТА № 1}

\noindent 1.\; Сформировать случайным образом связный ациклический ориентированный граф в соответствии с \textit{нормальным распределением} (параметры распределения задаются как константы, подобранные экспериментально) с вводимым пользователем количеством вершин. Распределение взято из справочника Вадзинского~\cite{vadzinsky}. Генерация происходит по степеням вершин.

\noindent 2.\; Посчитать экцентриситеты вершин, выделить центр графа и диаметральные вершины.

\noindent 3.\; Реализовать метод Шимбелла для сгенерированной с помощью заданного распределения весовой матрицы (пользователь вводит количество рёбер). В ответе представить минимальный и/или максимальный путь в виде матрицы. Предусмотреть генерацию только положительных, только отрицательных и смешанных весов на выбор пользователя.

\noindent 4.\; Определить возможность построения маршрута от одной заданной точки до другой (вершины вводит пользователь) и указывать количество таковых маршрутов.

\noindent\textit{Вариант: нормальное распределение (Normal distribution).}

\medskip
\noindent\textbf{ЛАБОРАТОРНАЯ РАБОТА № 2}

\noindent 1.\; Для заданных графов (случайно сгенерированных в предыдущей работе) выполнить \textit{обход вершин графа поиском в глубину}.

\noindent 2.\; Сгенерировать весовую матрицу (положительную или с отрицательными весами~--- выбирает пользователь) и найти кратчайший путь для двух выбранных пользователем вершин с помощью \textit{алгоритма Флойда-Уоршалла}. В результате выводится матрица расстояний, а также сам путь в виде последовательности вершин.

\noindent 3.\; Сравнить скорости работы реализованных алгоритмов (по количеству итераций).

\noindent\textit{Вариант: (с) обход вершин графа поиском в глубину; (j) алгоритм Флойда-Уоршалла.}

\medskip
\noindent\textbf{ЛАБОРАТОРНАЯ РАБОТА № 3}

\noindent 1.\; На основе сгенерированного ориентированного графа получить случайным образом матрицы пропускных способностей и стоимости.

\noindent 2.\; Для полученного графа найти максимальный поток по \textit{алгоритму Форда-Фалкерсона}.

\noindent 3.\; Вычислить заданный поток минимальной стоимости (в качестве величины потока брать значение, равное $\lfloor 2/3 \cdot \mathrm{max} \rfloor$, где $\mathrm{max}$~--- максимальный поток). Использовать ранее реализованный алгоритм Флойда-Уоршалла.

\medskip
\noindent\textbf{ЛАБОРАТОРНАЯ РАБОТА № 4}

\noindent 1.\; Для заданных неориентированных графов (случайно сгенерированных в первой работе) найти число остовных деревьев, используя матричную теорему Кирхгофа.

\noindent 2.\; Построить минимальный по весу остов для сгенерированного (неориентированного) взвешенного графа, используя \textit{алгоритм Борувки}. Полученный остов закодировать с помощью кода Прюфера и декодировать его (с сохранением весов).

\noindent 3.\; Реализовать на исходном графе и полученном остове (по выбору пользователя) \textit{нахождение минимального рёберного покрытия}.

\noindent\textit{Вариант: (b) алгоритм Борувки; (g) нахождение минимального рёберного покрытия.}

\medskip
\noindent\textbf{ЛАБОРАТОРНАЯ РАБОТА № 5}

\noindent 1.\; Для заданных неориентированных графов (случайно сгенерированных в первой работе) проверить, является ли граф эйлеровым. Если нет~--- модифицировать граф (логировать, что изменено). Построить эйлеров цикл.

\noindent 2.\; Построить кратчайший остов (алгоритмом Борувки из четвёртой работы). На основе данного остова получить \textit{фундаментальную систему разрезов} (вывести на экран); реализовать получение всех остальных разрезов на основе фундаментальных с помощью операции симметрической разности (выбранные разрезы вводит пользователь).

\noindent\textit{Вариант: 36, чётный~--- фундаментальная система разрезов.}

\medskip
Лабораторные работы выполнены на языке C++ в среде разработки CLion.

\newpage
	\section{Математическое описание}
	\subsection{Граф}

Простой граф $G(V, E)$ есть совокупность двух множеств~--- непустого множества $V$ и множества $E$ неупорядоченных пар различных элементов множества $V$. Множество $V$ называется множеством вершин, множество $E$ называется множеством рёбер.
\[
G(V,E) = \langle V, E\rangle,\quad V \neq \emptyset,\quad E \subseteq V \times V,\quad \{v,v\} \notin E,\quad v \in V,
\]
то есть множество $E$ состоит из 2-элементных подмножеств множества $V$.

Сопутствующие термины:
\begin{itemize}
  \item Вершина графа $G$ есть элемент множества вершин $v \in V(G)$;
  \item Ребро графа $G$ есть элемент множества рёбер $e \in E(G)$, или $e = \{v_1, v_2\}$, где $v_1 \in V(G)$, $v_2 \in V(G)$;
  \item Элементами графа $G$ называются его вершины $v \in V(G)$ и рёбра $e \in E(G)$ графа;
  \item Вершина $v$ инцидентна ребру $e$, если $v \in e$; тогда ещё говорят, что $e$ есть ребро при $v$;
  \item Соседние (смежные) вершины есть такие вершины $v_1$ и $v_2$, что $\{v_1,v_2\} \in E(G)$;
  \item Степень вершины $d(v)$ есть количество инцидентных ей рёбер.
\end{itemize}

	\subsection{Нормальное распределение}

Нормальное (гауссово) распределение~--- непрерывное распределение вероятностей. Плотность вероятности:
\[
f(x) = \frac{1}{\sigma\sqrt{2\pi}}\,\exp\!\left(-\frac{(x-\mu)^2}{2\sigma^2}\right),
\]
где $\mu$~--- математическое ожидание, $\sigma^2$~--- дисперсия.

Для генерации случайных степеней вершин используется приближённый метод из справочника Вадзинского~\cite{vadzinsky}. Случайная величина $X$, имеющая нормальное распределение $\mathcal{N}(\mu,\sigma^2)$, аппроксимируется следующим образом:
\[
X \approx \sqrt{\frac{12}{n}}\left(\sum_{i=1}^{n} r_i - \frac{n}{2}\right),
\]
где $r_i \sim U(0,1)$~--- независимые равномерно распределённые на $[0,1]$ случайные числа, $n$~--- количество слагаемых (параметр точности аппроксимации). Распределение взято из справочника по вероятностным распределениям~\cite{vadzinsky}.

	\subsection{Экцентриситет, центр и диаметральные вершины}

\textbf{Экцентриситетом} вершины $v$ в графе $G$ называется максимальное расстояние от $v$ до остальных вершин:
\[
\varepsilon(v) = \max_{u \in V}\, d(v, u),
\]
где $d(v, u)$~--- длина кратчайшего пути между вершинами $v$ и $u$.

\textbf{Радиусом} графа называется минимальный экцентриситет среди всех вершин:
\[
\mathrm{rad}(G) = \min_{v \in V}\, \varepsilon(v).
\]

\textbf{Диаметром} графа называется максимальный экцентриситет:
\[
\mathrm{diam}(G) = \max_{v \in V}\, \varepsilon(v).
\]

\textbf{Центром} графа называется множество вершин с минимальным экцентриситетом, равным радиусу:
\[
Z(G) = \{v \in V \mid \varepsilon(v) = \mathrm{rad}(G)\}.
\]

\textbf{Диаметральными} называются вершины с экцентриситетом, равным диаметру графа. Экцентриситеты вычисляются с помощью поиска в ширину (BFS) из каждой вершины за время $O(V(V+E))$.

	\subsection{Метод Шимбелла}

Метод Шимбелла позволяет находить минимальные (максимальные) пути между вершинами, состоящие из заданного количества рёбер.

\begin{itemize}
  \item Операция умножения двух величин $a$ и $b$ при возведении матрицы в степень соответствует их алгебраической сумме, то есть:
  \[a * b = b * a \;\to\; a + b = b + a,\qquad a * 0 = 0 * a = 0 \;\to\; a + 0 = 0.\]
  \item Операция сложения двух величин $a$ и $b$ заменяется выбором из этих величин максимума или минимума, то есть:
  \[a + b = b + a = \min(\max)\,\{a,\,b\}.\]
  Нули при этом игнорируются.
  \[w_{ij} = \min\bigl(\max(w_{i1}+w_{1j}),\ldots,(w_{in}+w_{nj})\bigr).\]
\end{itemize}

	\subsection{Обход графа поиском в глубину}

Обход в глубину~--- один из основных методов обхода графа, часто используемый для проверки связности, поиска цикла и компонент сильной связности и для топологической сортировки.

Общая идея алгоритма состоит в следующем: для каждой не пройденной вершины необходимо найти все не пройденные смежные вершины и повторить поиск для них.

\begin{enumerate}
  \item Выбираем любую вершину из ещё не пройденных, обозначим её как $u$;
  \item Запускаем процедуру $\mathrm{dfs}(u)$:
  \begin{itemize}
    \item Помечаем вершину $u$ как пройденную;
    \item Для каждой не пройденной смежной с $u$ вершины (назовём её $v$) запускаем $\mathrm{dfs}(v)$;
  \end{itemize}
  \item Повторяем шаги 1 и 2, пока все вершины не окажутся пройденными.
\end{enumerate}

Время работы алгоритма оценивается как $O(V+E)$.

	\subsection{Алгоритм Флойда-Уоршалла}

Алгоритм Флойда-Уоршалла предназначен для нахождения кратчайших путей между всеми парами вершин во взвешенном графе, в том числе с отрицательными весами рёбер (при отсутствии отрицательных циклов).

Обозначим через $d^{(k)}[i][j]$ длину кратчайшего пути от вершины $i$ до вершины $j$, проходящего только через промежуточные вершины из множества $\{1, 2, \ldots, k\}$. Рекуррентное соотношение:
\[
d^{(k)}[i][j] = \min\!\left(d^{(k-1)}[i][j],\; d^{(k-1)}[i][k] + d^{(k-1)}[k][j]\right).
\]

Начальные условия:
\[
d^{(0)}[i][j] = \begin{cases} w(i,j), & \text{если ребро } (i,j) \text{ существует},\\ \infty, & \text{иначе}. \end{cases}
\]

В результате $d^{(n)}[i][j]$ содержит длину кратчайшего пути от $i$ до $j$. В ответе выводится матрица расстояний, а также сам путь в виде последовательности вершин. Трудоёмкость алгоритма~--- $O(V^3)$.

	\subsection{Алгоритм Форда-Фалкерсона}

Алгоритм Форда-Фалкерсона решает задачу нахождения максимального потока в сети. Идея алгоритма заключается в следующем:

\begin{itemize}
  \item Изначально величине потока присваивается значение~0 для всех $u$ и $v$ в транспортной сети.
  \item Затем величина потока итеративно увеличивается посредством поиска увеличивающего пути от источника $s$ к стоку $t$.
  \item Процесс повторяется, пока можно найти увеличивающий путь.
\end{itemize}

В реализованном варианте Эдмондса-Карпа, где увеличивающие пути ищутся поиском в ширину, трудоёмкость алгоритма составляет $O(VE^2)$.

\medskip
\noindent\textbf{Алгоритм поиска заданного потока минимальной стоимости} может использовать любой из алгоритмов поиска кратчайшего пути из истока в сток. Для заданного суммарного потока~$n$ последовательно $n$~раз ищется кратчайший по стоимости путь потока~1 и далее на этом пути пропускная способность всех рёбер уменьшается на~1. Далее $n$ таких стоимостей суммируются. Целевая функция:
\[
\min\!\left(\sum_{(u,v)\in E} c(u,v)\cdot f(u,v)\right),
\]
где $E$~--- множество всех рёбер в графе, $c(u,v)$~--- стоимость потока по ребру $(u,v)$, $f(u,v)$~--- поток по ребру $(u,v)$. Асимптотика алгоритма~--- $O(n \cdot (V+E))$.

	\subsection{Остов}

Пусть $G = (V, E)$~--- граф. Остовный подграф графа $G = (V, E)$~--- подграф, содержащий все вершины. Остовный подграф, являющийся деревом, называется остовом. Несвязный граф не имеет остова. Связный граф может иметь много остовов.

Минимальное остовное дерево графа $G$~--- остовное дерево $T = (V, E_T)$ графа $G$, сумма весов рёбер которого минимальна.

	\subsection{Матричная теорема Кирхгофа}

\textbf{Матрица Кирхгофа.} Матрица $B$ графа $G$ размера $n \times n$ определяется следующим образом:
\[
B[i,j] = \begin{cases}
-1, & \text{если вершины с номерами } i \text{ и } j \text{ смежны};\\[2pt]
0,  & \text{если } i \neq j \text{ и вершины не смежны};\\[2pt]
\deg(v_i), & \text{если } i = j.
\end{cases}
\]

\textbf{Свойства матрицы Кирхгофа:}
\begin{enumerate}
  \item Суммы элементов в каждой строке и каждом столбце матрицы равны~0.
  \item Алгебраические дополнения всех элементов матрицы равны между собой.
  \item Определитель матрицы Кирхгофа равен нулю.
\end{enumerate}

\textbf{Матричная теорема Кирхгофа.} Число остовных деревьев в связном графе $G$ порядка $n \geq 2$ равно алгебраическому дополнению любого элемента матрицы Кирхгофа $B$ графа $G$.

	\subsection{Алгоритм Борувки}

Алгоритм Борувки~--- алгоритм построения минимального остовного дерева взвешенного связного неориентированного графа, предложенный Отакаром Борувкой в~1926~году.

Алгоритм работает следующим образом:
\begin{enumerate}
  \item Изначально каждая вершина образует отдельную компоненту (лес из одиночных вершин).
  \item Для каждой компоненты находится ребро минимального веса, соединяющее её с другой компонентой.
  \item Все найденные рёбра добавляются в остов, соответствующие компоненты объединяются.
  \item Шаги 2--3 повторяются, пока не останется одна компонента.
\end{enumerate}

На каждой итерации число компонент уменьшается как минимум вдвое, поэтому выполняется не более $O(\log V)$ итераций. Общая трудоёмкость~--- $O(E \log V)$.

	\subsection{Кодирование Прюфера}

Алгоритм кодирования Прюфера состоит из следующих шагов. Пока количество вершин больше двух:
\begin{itemize}
  \item Выбираем лист $v$ с минимальным номером.
  \item В код Прюфера добавляем номер вершины, смежной с $v$.
  \item Вершина $v$ и инцидентное ей ребро удаляются из дерева.
\end{itemize}

Полученная последовательность называется кодом Прюфера для заданного дерева.

Алгоритм декодирования: имеется массив кода Прюфера размера $n-2$ и массив всех вершин графа $[1\ldots n]$. Далее $n-2$ раза повторяется следующая процедура: берётся первый элемент массива кода Прюфера, и в массиве вершин производится поиск наименьшей вершины, не содержащейся в коде. Найденная вершина и текущий элемент кода составляют ребро дерева. Данные вершины удаляются из соответствующих массивов. В конце в массиве с вершинами остаётся две вершины~--- они составляют последнее ребро дерева. Асимптотика алгоритмов~--- $O(n)$.

	\subsection{Минимальное рёберное покрытие}

\textbf{Рёберным покрытием} графа $G = (V, E)$ называется такое подмножество рёбер $S \subseteq E$, что каждая вершина графа инцидентна хотя бы одному ребру из $S$. \textbf{Минимальным рёберным покрытием} называется рёберное покрытие наименьшей мощности.

Связь с максимальным паросочетанием выражается теоремой Галлаи: для графа без изолированных вершин
\[
\alpha'(G) + \beta'(G) = |V(G)|,
\]
где $\alpha'(G)$~--- размер максимального паросочетания, $\beta'(G)$~--- размер минимального рёберного покрытия.

Алгоритм нахождения минимального рёберного покрытия:
\begin{enumerate}
  \item Найти максимальное паросочетание $M$ в графе.
  \item Для каждой вершины, не покрытой паросочетанием $M$, добавить в покрытие произвольное инцидентное ей ребро.
\end{enumerate}

Трудоёмкость определяется сложностью нахождения максимального паросочетания.

	\subsection{Эйлеров граф}

Эйлеров граф~--- это граф, в котором существует эйлеров цикл, то есть простой цикл, проходящий по всем рёбрам графа.

Граф является эйлеровым тогда и только тогда, когда он связен и степени всех его вершин чётны.

Эйлеровым путём называется простой путь, содержащий все рёбра. В данной работе производится поиск эйлерова цикла. Сначала осуществляется проверка, является ли граф эйлеровым. Если нет, то граф достраивается так, чтобы степень каждой вершины была чётной.

Для поиска эйлерова цикла используется алгоритм, основанный на обходе в глубину (DFS):
\begin{itemize}
  \item Выбрать произвольную вершину графа и поместить её в стек.
  \item Пока стек не пуст, повторять следующие действия:
  \begin{itemize}
    \item извлечь вершину из стека;
    \item если у этой вершины есть непосещённые соседние вершины, выбрать одну из них и поместить её в стек;
    \item если у вершины нет непосещённых соседних вершин, добавить её в список эйлерова цикла.
  \end{itemize}
  \item Вернуться к списку эйлерова цикла в обратном порядке.
\end{itemize}

Асимптотика алгоритма поиска эйлерова цикла оценивается как $O(V+E)$.

	\subsection{Фундаментальная система разрезов}

\textbf{Разрезом} (разрезающим множеством) в графе $G = (V, E)$ называется такое множество рёбер $C \subseteq E$, удаление которых увеличивает число компонент связности графа.

\textbf{Фундаментальным разрезом} относительно остовного дерева $T$ называется разрез, соответствующий одному ребру дерева $t \in E_T$: при удалении $t$ дерево $T$ распадается на две компоненты $T_1$ и $T_2$; фундаментальный разрез $C_t$~--- множество всех рёбер исходного графа $G$, у которых одна конечная вершина лежит в $T_1$, а другая в~$T_2$.

\textbf{Фундаментальная система разрезов} строится для выбранного остовного дерева $T$ и содержит ровно $|V|-1$ разрезов (по числу рёбер дерева).

Все остальные разрезы графа можно получить с помощью операции симметрической разности фундаментальных разрезов:
\[
C = C_{t_1} \mathbin{\triangle} C_{t_2} \mathbin{\triangle} \cdots \mathbin{\triangle} C_{t_k},
\]
где $\triangle$~--- операция симметрической разности множеств.

\clearpage
\section{Особенности реализации}

\subsection{Структура проекта}

Программа реализована на языке C++17 и имеет модульную структуру. Общие структуры данных и функции генерации вынесены в каталог \texttt{shared}, алгоритмы разделены по лабораторным работам, а каталог \texttt{lab\_combined} содержит объединяющую консольную программу с меню. Полный исходный код приведён в приложениях А--Е.

\begin{lstlisting}[caption={Структура проекта}, label={lst:project-structure}]
graph-theory-labs/
|-- CMakeLists.txt
|-- README.md
|-- shared/
|   |-- constants.h
|   |-- graph.h
|   |-- graph.cpp
|   |-- distribution.h
|   |-- distribution.cpp
|   |-- dag_generator.h
|   |-- dag_generator.cpp
|   |-- weight_matrix.h
|   `-- weight_matrix.cpp
|-- lab1/
|   |-- eccentricity.h
|   |-- eccentricity.cpp
|   |-- shimbell.h
|   |-- shimbell.cpp
|   |-- path_counter.h
|   `-- path_counter.cpp
|-- lab2/
|   |-- dfs.h
|   |-- dfs.cpp
|   |-- floyd_warshall.h
|   `-- floyd_warshall.cpp
|-- lab3/
|   |-- ford_fulkerson.h
|   |-- ford_fulkerson.cpp
|   |-- min_cost_flow.h
|   `-- min_cost_flow.cpp
|-- lab4/
|   |-- kirchhoff.h
|   |-- kirchhoff.cpp
|   |-- boruvka.h
|   |-- boruvka.cpp
|   |-- edge_cover.h
|   |-- edge_cover.cpp
|   |-- prufer.h
|   `-- prufer.cpp
|-- lab5/
|   |-- eulerian.h
|   |-- eulerian.cpp
|   |-- fundamental_cutsets.h
|   `-- fundamental_cutsets.cpp
`-- lab_combined/
    |-- CMakeLists.txt
    `-- main.cpp
\end{lstlisting}

Каталог \texttt{shared} содержит класс графа, генератор нормального распределения, генератор графа и генератор весовой матрицы. Каталоги \texttt{lab1}--\texttt{lab5} содержат алгоритмы соответствующих лабораторных работ. Файл \texttt{lab\_combined/main.cpp} связывает все модули в единое меню пользователя.

\subsection{Класс Graph}

Класс \texttt{Graph} является базовой структурой для хранения ориентированного или неориентированного графа. Его объявление приведено в листинге~\ref{lst:shared-graph-h}, а реализация методов -- в листинге~\ref{lst:shared-graph-cpp}. В классе хранятся количество вершин \texttt{n}, признак ориентированности \texttt{directed} и матрица смежности \texttt{adj}.

Функция \texttt{Graph(int n, bool directed)} создаёт пустой граф нужного типа.

Вход: количество вершин \texttt{n}, логический признак ориентированности \texttt{directed}.

Выход: объект графа с матрицей смежности размера $n \times n$, заполненной нулями.

Функция \texttt{addEdge} добавляет ребро или дугу между двумя вершинами. Если граф неориентированный, то матрица смежности заполняется симметрично.

Вход: номера вершин \texttt{u} и \texttt{v}.

Выход: изменённая матрица смежности графа.

Функция \texttt{hasEdge} проверяет существование ребра или дуги.

Вход: номера вершин \texttt{u} и \texttt{v}.

Выход: значение \texttt{true}, если ребро существует, иначе \texttt{false}.

Функция \texttt{neighbors} возвращает все вершины, смежные с заданной. Для ориентированного графа возвращаются исходящие соседи, для неориентированного -- все соседние вершины.

Вход: номер вершины \texttt{u}.

Выход: вектор номеров соседних вершин.

Функции \texttt{printAdjMatrix} и \texttt{printFull} используются для вывода графа в консоль. Первая выводит только матрицу смежности, а вторая также печатает список смежности, список рёбер и, при наличии, весовую матрицу.

Вход: объект графа и, для \texttt{printFull}, необязательная весовая матрица.

Выход: текстовое представление графа в консоли.

\subsection{Генерация распределения, графа и весовой матрицы}

Функции генерации распределения описаны в листингах~\ref{lst:shared-distribution-h} и~\ref{lst:shared-distribution-cpp}. Для генерации используется приближённая формула нормального распределения через сумму равномерно распределённых случайных величин.

Функция \texttt{normalSample} генерирует одно случайное число, приближённо распределённое по нормальному закону.

Вход: количество равномерных случайных величин \texttt{n}, используемых в аппроксимации.

Выход: вещественное число, соответствующее приближённому нормальному распределению.

Функция \texttt{sampleDegree} переводит значение нормального распределения в допустимую степень вершины.

Вход: максимальная допустимая степень \texttt{maxDeg}.

Выход: целое число в диапазоне от 1 до \texttt{maxDeg}.

Функция \texttt{generateDegreeSequence} строит последовательность степеней для заданного количества вершин. Сумма степеней корректируется до значения $2(n-1)$, что соответствует связному ациклическому неориентированному графу.

Вход: количество вершин \texttt{numVertices}.

Выход: вектор степеней вершин.

Генерация графа реализована в файлах \texttt{dag\_generator.h} и \texttt{dag\_generator.cpp}, приведённых в листингах~\ref{lst:shared-dag-generator-h} и~\ref{lst:shared-dag-generator-cpp}. Основная функция \texttt{generateDAG} создаёт ориентированный ациклический граф или неориентированный граф, полученный на основе ориентированной модели.

Вход: количество вершин \texttt{n} и признак ориентированности \texttt{directed}.

Выход: объект \texttt{Graph} со сгенерированной матрицей смежности.

Для ориентированного графа сначала строится случайный топологический порядок вершин, затем добавляется путь, обеспечивающий связность, после чего добавляются дополнительные дуги только вперёд по топологическому порядку. Благодаря этому граф остаётся ациклическим. Для неориентированного графа направления дуг игнорируются, и граф становится симметричным.

Генерация весовой матрицы реализована в листингах~\ref{lst:shared-weight-matrix-h} и~\ref{lst:shared-weight-matrix-cpp}. Функция \texttt{generateWeightMatrix} заполняет веса только для существующих рёбер графа.

Вход: граф \texttt{g} и режим генерации весов \texttt{mode}.

Выход: матрица весов, где для отсутствующих рёбер используется значение \texttt{INF}.

Режим \texttt{POSITIVE} создаёт только положительные веса, \texttt{NEGATIVE} -- только отрицательные, а \texttt{MIXED} -- смешанные веса. Для неориентированного графа веса записываются симметрично.

\subsection{Лабораторная работа №1}

В первой лабораторной работе реализованы вычисление метрических характеристик графа, метод Шимбелла и поиск маршрутов. Код расположен в файлах \texttt{eccentricity.cpp}, \texttt{shimbell.cpp} и \texttt{path\_counter.cpp}; полные тексты приведены в листингах~\ref{lst:lab1-eccentricity-cpp},~\ref{lst:lab1-shimbell-cpp} и~\ref{lst:lab1-path-counter-cpp}.

Функция \texttt{bfs} выполняет поиск в ширину из заданной вершины и используется для вычисления расстояний до остальных вершин.

Вход: граф \texttt{g}, начальная вершина \texttt{src}.

Выход: вектор расстояний от вершины \texttt{src}; значение $-1$ означает, что вершина недостижима.

Функция \texttt{computeMetrics} вычисляет эксцентриситеты, радиус, диаметр, центр и диаметральные вершины графа. Для каждой вершины запускается \texttt{bfs}, после чего из полученных расстояний выбирается максимальное достижимое расстояние.

Вход: граф \texttt{g}.

Выход: структура \texttt{GraphMetrics}, содержащая матрицу расстояний, эксцентриситеты, радиус, диаметр, центр и диаметральные вершины.

Функция \texttt{shimbell} реализует метод Шимбелла для путей, состоящих ровно из $k$ рёбер. Для минимальных путей используется тропическое умножение с выбором минимума, а для максимальных -- с выбором максимума.

Вход: весовая матрица \texttt{W}, количество рёбер \texttt{k}.

Выход: пара матриц: матрица минимальных путей и матрица максимальных путей.

Функция \texttt{findAllPaths} ищет все простые маршруты между двумя вершинами. Реализация использует рекурсивный поиск в глубину с массивом посещённых вершин, чтобы не допускать повторения вершин в одном маршруте.

Вход: граф \texttt{g}, начальная вершина \texttt{src}, конечная вершина \texttt{dst}.

Выход: вектор маршрутов, где каждый маршрут представлен последовательностью вершин.

Функции \texttt{pathExists} и \texttt{countPaths} являются вспомогательными. Первая проверяет наличие хотя бы одного маршрута, вторая возвращает количество найденных маршрутов. Функция \texttt{printPaths} выводит маршруты в консоль.

Вход: граф или уже найденный список маршрутов.

Выход: логическое значение, количество маршрутов или текстовый вывод в консоль.

\subsection{Лабораторная работа №2}

Во второй лабораторной работе реализованы обход графа поиском в глубину и алгоритм Флойда-Уоршалла. Код приведён в листингах~\ref{lst:lab2-dfs-cpp} и~\ref{lst:lab2-floyd-warshall-cpp}.

Функция \texttt{dfs} выполняет обход графа в глубину из заданной вершины. В реализации используется стек, поэтому алгоритм работает итеративно. Дополнительно подсчитывается количество итераций.

Вход: граф \texttt{g}, стартовая вершина \texttt{start}.

Выход: порядок обхода, список недостижимых вершин и количество итераций, выведенные в консоль.

Функция \texttt{floydWarshall} находит кратчайшие пути между всеми парами вершин. Матрица \texttt{T} хранит расстояния, а матрица \texttt{H} хранит информацию для восстановления пути.

Вход: граф \texttt{g}, весовая матрица \texttt{W}.

Выход: структура \texttt{FloydResult}, содержащая матрицу расстояний, матрицу восстановления пути, количество итераций и признак отрицательного цикла.

Функция \texttt{printFloydResult} выводит результат работы алгоритма и восстанавливает путь между двумя выбранными пользователем вершинами. Если путь не существует, программа выводит соответствующее сообщение.

Вход: результат алгоритма Флойда-Уоршалла, исходная весовая матрица, количество вершин, начальная и конечная вершины.

Выход: матрицы \texttt{C}, \texttt{T}, \texttt{H}, кратчайший путь и его длина в консоли.

\subsection{Лабораторная работа №3}

В третьей лабораторной работе реализованы генерация матриц пропускных способностей и стоимостей, алгоритм Форда-Фалкерсона и поиск потока минимальной стоимости. Код расположен в листингах~\ref{lst:lab3-ford-fulkerson-cpp} и~\ref{lst:lab3-min-cost-flow-cpp}.

Функция \texttt{generateCapacityMatrix} создаёт матрицу пропускных способностей для существующих дуг ориентированного графа.

Вход: ориентированный граф \texttt{g}.

Выход: матрица пропускных способностей, где ноль означает отсутствие дуги.

Функция \texttt{generateCostMatrix} создаёт матрицу стоимости прохождения единицы потока по каждой дуге.

Вход: ориентированный граф \texttt{g}.

Выход: матрица стоимостей, соответствующая структуре графа.

Функция \texttt{fordFulkerson} реализует алгоритм Форда-Фалкерсона в варианте Эдмондса-Карпа, то есть увеличивающий путь ищется с помощью поиска в ширину. На каждом шаге находится пропускная способность пути, после чего остаточная сеть обновляется.

Вход: количество вершин, матрица пропускных способностей, исток \texttt{src}, сток \texttt{dst}.

Выход: структура \texttt{FFResult}, содержащая матрицу потока и значение максимального потока.

Функция \texttt{fordFulkersonSilent} выполняет тот же алгоритм, но без вывода промежуточных путей. Она используется внутри алгоритма потока минимальной стоимости.

Вход: количество вершин, матрица пропускных способностей, исток и сток.

Выход: значение максимального потока и матрица потока без печати в консоль.

Функция \texttt{minCostFlow} вычисляет поток минимальной стоимости для величины $\lfloor 2/3 \cdot \varphi_{\max} \rfloor$. Сначала находится максимальный поток, затем в остаточной сети последовательно выбираются кратчайшие по стоимости пути, найденные алгоритмом Флойда-Уоршалла.

Вход: количество вершин, матрица пропускных способностей, матрица стоимостей, исток и сток.

Выход: структура \texttt{MCFResult}, содержащая максимальный поток, требуемый поток, достигнутый поток, итоговую стоимость и матрицу потока.

Функция \texttt{printMCFResult} выводит итоговый поток минимальной стоимости и проверочный расчёт стоимости по дугам.

Вход: результат потока минимальной стоимости, матрица пропускных способностей, матрица стоимостей.

Выход: таблица дуг с положительным потоком и итоговая стоимость.

\subsection{Лабораторная работа №4}

В четвёртой лабораторной работе реализованы матричная теорема Кирхгофа, алгоритм Борувки, минимальное рёберное покрытие и код Прюфера. Полный код приведён в листингах~\ref{lst:lab4-kirchhoff-cpp},~\ref{lst:lab4-boruvka-cpp},~\ref{lst:lab4-edge-cover-cpp} и~\ref{lst:lab4-prufer-cpp}.

Функция \texttt{countSpanningTrees} считает количество остовных деревьев графа. Для этого строится матрица Кирхгофа, удаляется первая строка и первый столбец, после чего вычисляется определитель полученной матрицы.

Вход: неориентированный граф \texttt{g}.

Выход: количество остовных деревьев графа.

Функция \texttt{boruvka} строит минимальное остовное дерево. На каждой итерации для каждой компоненты выбирается минимальное исходящее ребро, после чего компоненты объединяются.

Вход: неориентированный граф \texttt{g}, весовая матрица \texttt{W}.

Выход: структура \texttt{BoruvkaResult}, содержащая рёбра минимального остова и его суммарный вес.

Функция \texttt{minEdgeCover} находит минимальное рёберное покрытие. Сначала строится максимальное паросочетание, затем для непокрытых вершин добавляются инцидентные рёбра.

Вход: граф \texttt{g}, весовая матрица, признак работы с остовом или исходным графом, список рёбер остова.

Выход: структура \texttt{EdgeCoverResult}, содержащая рёбра покрытия и размер максимального паросочетания.

Функция \texttt{pruferEncode} кодирует минимальный остов кодом Прюфера. На каждом шаге удаляется лист с минимальным номером, а в код записывается его сосед. Вес удалённого ребра сохраняется отдельно.

Вход: список рёбер остова, количество вершин.

Выход: код Прюфера и список весов рёбер в порядке удаления.

Функция \texttt{pruferDecode} восстанавливает дерево по коду Прюфера и сохранённым весам. На каждом шаге выбирается минимальный лист, отсутствующий в текущем коде.

Вход: код Прюфера, список весов, количество вершин.

Выход: список рёбер восстановленного дерева с весами.

\subsection{Лабораторная работа №5}

В пятой лабораторной работе реализованы проверка графа на эйлеровость, построение эйлерова цикла и фундаментальная система разрезов. Код приведён в листингах~\ref{lst:lab5-eulerian-cpp} и~\ref{lst:lab5-fundamental-cutsets-cpp}.

Функция \texttt{buildEulerianCycle} проверяет связность графа и чётность степеней всех вершин. Если граф не является эйлеровым, программа пытается изменить его, добавляя или удаляя рёбра так, чтобы степени стали чётными и связность сохранилась.

Вход: неориентированный граф \texttt{g}.

Выход: структура \texttt{EulerianResult}, содержащая исходные степени, список нечётных вершин, список изменений и построенный эйлеров цикл.

Функция \texttt{hierholzer} строит эйлеров цикл после того, как граф стал эйлеровым. Алгоритм последовательно удаляет пройденные рёбра и формирует цикл в обратном порядке.

Вход: матрица смежности эйлерова графа.

Выход: последовательность вершин эйлерова цикла.

Функция \texttt{buildFundamentalCutsets} строит фундаментальную систему разрезов относительно минимального остова. Для каждого ребра остова оно временно удаляется, после чего находятся две компоненты дерева и все рёбра исходного графа, соединяющие эти компоненты.

Вход: исходный граф \texttt{g}, список рёбер минимального остова.

Выход: вектор фундаментальных разрезов.

Функция \texttt{printCutsetSymmetricDifference} вычисляет симметрическую разность выбранных фундаментальных разрезов. Ребро, встречающееся два раза, удаляется из результата, а ребро, встречающееся один раз, остаётся.

Вход: список фундаментальных разрезов и номера выбранных разрезов.

Выход: множество рёбер, полученное операцией симметрической разности.

\subsection{Функция main и меню программы}

Объединение всех лабораторных работ выполнено в файле \texttt{lab\_combined/main.cpp}, полный код которого приведён в листинге~\ref{lst:lab-combined-main-cpp}. В программе используются глобальные переменные для хранения текущего графа, весовой матрицы, матриц пропускных способностей и стоимостей, минимального остова и результата кодирования Прюфера.

Функция \texttt{readInt} обеспечивает безопасный ввод целого числа. Если пользователь вводит некорректные данные, поток ввода очищается, и запрос повторяется.

Вход: текст приглашения \texttt{prompt}.

Выход: корректно считанное целое число.

Функции \texttt{requireGraph}, \texttt{requireWeight}, \texttt{requireFlowMatrices} и \texttt{requireMST} проверяют, были ли подготовлены необходимые данные перед запуском алгоритма.

Вход: текущее состояние программы.

Выход: сообщение об ошибке, если нужные данные отсутствуют.

Функции меню, например \texttt{menuGenerateGraph}, \texttt{menuShimbell}, \texttt{menuFloydWarshall}, \texttt{menuFordFulkerson}, \texttt{menuBoruvka} и другие, считывают параметры пользователя, проверяют корректность состояния программы и вызывают соответствующие алгоритмы.

Вход: выбор пользователя и дополнительные значения, введённые в консоль.

Выход: запуск выбранного алгоритма и вывод результата в консоль.

Функция \texttt{main} организует главный цикл программы. Пока пользователь не выберет пункт выхода, на экран выводится меню, считывается номер команды и вызывается соответствующая функция обработки.

Вход: команды пользователя из консольного меню.

Выход: последовательное выполнение лабораторных алгоритмов до завершения программы.

Таким образом, объединяющая программа позволяет выполнять лабораторные работы в одном интерфейсе: сначала генерируется граф, затем строятся необходимые матрицы, после чего пользователь может запускать алгоритмы из разных разделов работы.

\subsection{Подробное описание функций}

Ниже приведено более подробное описание основных функций программы. Для каждой функции указано её полное имя так, как оно используется в исходном коде, затем описаны назначение, входные данные и результат работы. Сам исходный код не дублируется в данном разделе, потому что он полностью приведён в приложениях А--Е; здесь код используется как основа для объяснения реализации. Названия функций в данном подразделе являются гиперссылками на соответствующие листинги с исходным кодом.

\subsubsection*{Общие структуры данных и генерация}

\noindent\textbf{Функция} \coderef{Graph::Graph()}{lst:shared-graph-h}\\
Конструктор по умолчанию создаёт пустой объект графа. Он нужен для глобальных переменных и ситуаций, когда граф будет сгенерирован позже через меню программы. Внутри конструктора количество вершин устанавливается равным нулю, а граф считается неориентированным до момента явного создания нового объекта.

\textbf{Вход:} входные параметры отсутствуют.

\textbf{Выход:} пустой объект \texttt{Graph}, для которого метод \coderef{isEmpty()}{lst:shared-graph-h} возвращает значение \texttt{true}.

\noindent\textbf{Функция} \coderef{Graph::Graph(int n, bool directed)}{lst:shared-graph-cpp}\\
Основной конструктор класса \texttt{Graph}. Он создаёт матрицу смежности размера $n \times n$ и заполняет её нулями. Параметр \texttt{directed} определяет, будет ли граф ориентированным. Благодаря хранению графа в виде матрицы смежности проверка наличия ребра выполняется за постоянное время.

\textbf{Вход:} количество вершин \texttt{n}; логический параметр \texttt{directed}, равный \texttt{true} для ориентированного графа и \texttt{false} для неориентированного.

\textbf{Выход:} объект графа с подготовленной матрицей смежности и сохранённым типом графа.

\noindent\textbf{Функция} \coderef{Graph::addEdge(int u, int v)}{lst:shared-graph-cpp}\\
Метод добавляет ребро между вершинами \texttt{u} и \texttt{v}. Если граф ориентированный, то в матрице смежности устанавливается только элемент \texttt{adj[u][v]}. Если граф неориентированный, дополнительно устанавливается симметричный элемент \texttt{adj[v][u]}, поэтому одно ребро хранится сразу в двух ячейках матрицы.

\textbf{Вход:} номера вершин \texttt{u} и \texttt{v} во внутренней нумерации, начиная с нуля.

\textbf{Выход:} изменённая матрица смежности текущего объекта \texttt{Graph}.

\noindent\textbf{Функция} \coderef{Graph::hasEdge(int u, int v) const}{lst:shared-graph-cpp}\\
Метод проверяет, существует ли ребро или дуга из вершины \texttt{u} в вершину \texttt{v}. Реализация обращается к одной ячейке матрицы смежности и сравнивает её с единицей.

\textbf{Вход:} номера вершин \texttt{u} и \texttt{v}.

\textbf{Выход:} значение \texttt{true}, если \texttt{adj[u][v] == 1}; иначе значение \texttt{false}.

\noindent\textbf{Функция} \coderef{Graph::neighbors(int u) const}{lst:shared-graph-cpp}\\
Метод формирует список всех вершин, достижимых из вершины \texttt{u} за один шаг. Для ориентированного графа это исходящие соседи, а для неориентированного графа это все смежные вершины, потому что матрица смежности симметрична.

\textbf{Вход:} номер вершины \texttt{u}.

\textbf{Выход:} вектор \lstinline|std::vector<int>| с номерами соседних вершин.

\noindent\textbf{Функция} \coderef{Graph::printAdjMatrix() const}{lst:shared-graph-cpp}\\
Метод выводит матрицу смежности в консоль. Вершины при выводе переводятся из внутренней нумерации $0,1,\ldots,n-1$ в пользовательскую нумерацию $1,2,\ldots,n$. Для отсутствующего ребра выводится \texttt{inf}, а для существующего ребра -- значение \texttt{1}.

\textbf{Вход:} текущий объект \texttt{Graph}.

\textbf{Выход:} таблица матрицы смежности в консоли.

\noindent\textbf{Функция} \coderef{Graph::printFull(const std::vector<std::vector<double>>* weightMatrix) const}{lst:shared-graph-cpp}\\
Метод выводит расширенное представление графа: список смежности, список рёбер, матрицу смежности и, если она была передана, весовую матрицу. Указатель на весовую матрицу сделан необязательным, поэтому одна функция используется и до генерации весов, и после неё.

\textbf{Вход:} текущий граф; необязательный указатель \texttt{weightMatrix} на матрицу весов.

\textbf{Выход:} полное текстовое описание графа в консоли.

\noindent\textbf{Функция} \coderef{Graph::isEmpty() const}{lst:shared-graph-h}\\
Метод используется для быстрой проверки, был ли уже создан граф. Он возвращает истину, если число вершин равно нулю.

\textbf{Вход:} текущий объект \texttt{Graph}.

\textbf{Выход:} логическое значение, показывающее, пустой граф или нет.

\noindent\textbf{Функция} \coderef{double normalSample(int n)}{lst:shared-distribution-cpp}\\
Функция генерирует одно значение, приближённо распределённое по нормальному закону. Для этого используется сумма \texttt{n} равномерно распределённых случайных величин и нормировка по формуле, указанной в комментарии исходного файла. Такой подход удобен, потому что он прост в реализации и не требует отдельного сложного генератора нормального распределения.

\textbf{Вход:} число \texttt{n}, задающее количество равномерных случайных величин для аппроксимации.

\textbf{Выход:} вещественное число типа \texttt{double}, имитирующее выборку из нормального распределения.

\noindent\textbf{Функция} \coderef{int sampleDegree(int maxDeg)}{lst:shared-distribution-cpp}\\
Функция переводит случайное значение нормального распределения в целую степень вершины. Полученное число ограничивается диапазоном от \texttt{1} до \texttt{maxDeg}, чтобы каждая вершина получила допустимую положительную степень.

\textbf{Вход:} максимальная допустимая степень \texttt{maxDeg}.

\textbf{Выход:} целое число, представляющее степень вершины.

\noindent\textbf{Функция} \coderef{std::vector<int> generateDegreeSequence(int numVertices)}{lst:shared-distribution-cpp}\\
Функция создаёт последовательность степеней для заданного количества вершин. Сначала для каждой вершины генерируется степень, затем сумма степеней корректируется так, чтобы она соответствовала требуемой модели генерации графа. Эта последовательность далее используется при построении ориентированного ациклического графа.

\textbf{Вход:} количество вершин \texttt{numVertices}.

\textbf{Выход:} вектор целых чисел, где каждый элемент соответствует степени одной вершины.

\noindent\textbf{Функция} \coderef{Graph generateDAG(int n, bool directed)}{lst:shared-dag-generator-cpp}\\
Функция является основной точкой генерации графа. Если параметр \texttt{directed} равен \texttt{true}, строится ориентированный ациклический граф. Для этого создаётся случайный топологический порядок, добавляется путь для связности, а дополнительные дуги проводятся только вперёд по этому порядку. Если параметр \texttt{directed} равен \texttt{false}, сначала строится ориентированная модель, а затем направления рёбер игнорируются и получается неориентированный граф.

\textbf{Вход:} количество вершин \texttt{n}; признак ориентированности \texttt{directed}.

\textbf{Выход:} сгенерированный объект \texttt{Graph}.

\noindent\textbf{Функция} \coderef{Matrix generateWeightMatrix(const Graph& g, WeightMode mode)}{lst:shared-weight-matrix-cpp}\\
Функция создаёт матрицу весов для уже сгенерированного графа. Вес назначается только там, где в графе существует ребро. Для отсутствующих рёбер используется специальное значение бесконечности, что важно для алгоритма Шимбелла и алгоритма Флойда-Уоршалла. Режим \texttt{POSITIVE} создаёт положительные веса, \texttt{NEGATIVE} -- отрицательные, а \texttt{MIXED} -- смешанные.

\textbf{Вход:} граф \texttt{g}; режим генерации весов \texttt{mode}.

\textbf{Выход:} матрица \texttt{Matrix} размера $n \times n$ с весами существующих рёбер.

\subsubsection*{Функции лабораторной работы №1}

\noindent\textbf{Функция} \coderef{std::vector<int> bfs(const Graph& g, int src)}{lst:lab1-eccentricity-cpp}\\
Функция выполняет поиск в ширину из вершины \texttt{src}. Она используется как вспомогательная часть вычисления метрических характеристик графа. Поиск идёт по списку соседей, полученному из матрицы смежности. Если вершина достижима, для неё записывается кратчайшее расстояние в количестве рёбер; если вершина недостижима, значение остаётся равным \texttt{-1}.

\textbf{Вход:} граф \texttt{g}; начальная вершина \texttt{src}.

\textbf{Выход:} вектор расстояний от вершины \texttt{src} до всех остальных вершин.

\noindent\textbf{Функция} \coderef{GraphMetrics computeMetrics(const Graph& g)}{lst:lab1-eccentricity-cpp}\\
Функция вычисляет основные метрические характеристики графа. Для каждой вершины вызывается \coderef{bfs}{lst:lab1-eccentricity-cpp}, после чего формируется матрица расстояний. Эксцентриситет вершины определяется как максимальное расстояние от неё до достижимых вершин. Затем среди эксцентриситетов выбираются радиус, диаметр, центральные вершины и диаметральные вершины.

\textbf{Вход:} граф \texttt{g}.

\textbf{Выход:} структура \texttt{GraphMetrics}, содержащая матрицу расстояний, эксцентриситеты, радиус, диаметр, центр и диаметральные вершины.

\noindent\textbf{Функция} \coderef{void printMetrics(const GraphMetrics& m)}{lst:lab1-eccentricity-cpp}\\
Функция отвечает только за вывод результатов, полученных в \coderef{computeMetrics}{lst:lab1-eccentricity-cpp}. Она печатает матрицу расстояний, список эксцентриситетов, радиус, диаметр, центр и диаметральные вершины в удобном для проверки виде.

\textbf{Вход:} структура \texttt{GraphMetrics}.

\textbf{Выход:} форматированный вывод метрических характеристик в консоли.

\noindent\textbf{Функция} \coderef{std::pair<Matrix, Matrix> shimbell(const Matrix& W, int k)}{lst:lab1-shimbell-cpp}\\
Функция реализует метод Шимбелла для нахождения путей, состоящих ровно из \texttt{k} рёбер. Внутри используются операции тропического умножения матриц: для минимальных путей выбирается минимум суммы весов, а для максимальных путей -- максимум. При \texttt{k = 0} возвращается аналог единичной матрицы, где путь из вершины в саму себя имеет длину ноль.

\textbf{Вход:} весовая матрица \texttt{W}; количество рёбер в пути \texttt{k}.

\textbf{Выход:} пара матриц: первая содержит минимальные веса путей, вторая -- максимальные веса путей для ровно \texttt{k} шагов.

\noindent\textbf{Функция} \coderef{void printMatrix(const Matrix& M)}{lst:lab1-shimbell-cpp}\\
Функция выводит матрицу весов на экран. Большие положительные значения отображаются как \texttt{INF}, большие отрицательные значения -- как \texttt{-INF}. Это делает вывод понятным, так как бесконечность означает отсутствие допустимого пути.

\textbf{Вход:} матрица \texttt{M}.

\textbf{Выход:} таблица значений матрицы в консоли.

\noindent\textbf{Функция} \coderef{std::vector<std::vector<int>> findAllPaths(const Graph& g, int src, int dst)}{lst:lab1-path-counter-cpp}\\
Функция находит все простые маршруты между вершинами \texttt{src} и \texttt{dst}. Для этого применяется рекурсивный поиск в глубину. Массив посещённых вершин не позволяет одной и той же вершине повторяться в одном маршруте, поэтому результат содержит именно простые пути.

\textbf{Вход:} граф \texttt{g}; начальная вершина \texttt{src}; конечная вершина \texttt{dst}.

\textbf{Выход:} список маршрутов, где каждый маршрут представлен вектором номеров вершин.

\noindent\textbf{Функция} \coderef{bool pathExists(const Graph& g, int src, int dst)}{lst:lab1-path-counter-cpp}\\
Функция является удобной обёрткой над поиском маршрутов. Она вызывает поиск путей и проверяет, найден ли хотя бы один путь между выбранными вершинами.

\textbf{Вход:} граф \texttt{g}; начальная вершина \texttt{src}; конечная вершина \texttt{dst}.

\textbf{Выход:} значение \texttt{true}, если путь существует; иначе \texttt{false}.

\noindent\textbf{Функция} \coderef{int countPaths(const Graph& g, int src, int dst)}{lst:lab1-path-counter-cpp}\\
Функция подсчитывает количество простых маршрутов между двумя вершинами. Она использует тот же результат, что и \coderef{findAllPaths}{lst:lab1-path-counter-cpp}, но возвращает только число маршрутов.

\textbf{Вход:} граф \texttt{g}; начальная вершина \texttt{src}; конечная вершина \texttt{dst}.

\textbf{Выход:} целое число, равное количеству найденных маршрутов.

\noindent\textbf{Функция} \coderef{void printPaths(const std::vector<std::vector<int>>& paths, int src, int dst)}{lst:lab1-path-counter-cpp}\\
Функция выводит найденные маршруты в пользовательской нумерации вершин. Если список пуст, она сообщает, что маршрутов между выбранными вершинами нет.

\textbf{Вход:} список маршрутов \texttt{paths}; начальная вершина \texttt{src}; конечная вершина \texttt{dst}.

\textbf{Выход:} текстовый список маршрутов в консоли.

\subsubsection*{Функции лабораторной работы №2}

\noindent\textbf{Функция} \coderef{void dfs(const Graph& g, int start)}{lst:lab2-dfs-cpp}\\
Функция выполняет обход графа поиском в глубину, начиная с вершины \texttt{start}. В реализации используется стек, поэтому алгоритм не зависит от глубины рекурсии. Во время работы фиксируется порядок посещения вершин и количество итераций. После обхода программа также показывает вершины, которые остались недостижимыми из выбранной начальной вершины.

\textbf{Вход:} граф \texttt{g}; стартовая вершина \texttt{start}.

\textbf{Выход:} порядок обхода, список недостижимых вершин и количество итераций в консоли.

\noindent\textbf{Функция} \coderef{FloydResult floydWarshall(const Graph& g, const std::vector<std::vector<double>>& W)}{lst:lab2-floyd-warshall-cpp}\\
Функция реализует алгоритм Флойда-Уоршалла для поиска кратчайших путей между всеми парами вершин. Матрица \texttt{T} хранит текущие расстояния, а матрица \texttt{H} используется для восстановления маршрута. После выполнения основного тройного цикла проверяются диагональные элементы матрицы расстояний, чтобы обнаружить отрицательные циклы.

\textbf{Вход:} граф \texttt{g}; весовая матрица \texttt{W}, в которой диагональ для этого алгоритма равна нулю.

\textbf{Выход:} структура \texttt{FloydResult} с матрицей расстояний, матрицей восстановления пути, числом итераций и признаком отрицательного цикла.

\noindent\textbf{Функция} \coderef{void printFloydResult(const FloydResult& res, const std::vector<std::vector<double>>& W, int n, int src, int dst)}{lst:lab2-floyd-warshall-cpp}\\
Функция выводит результат алгоритма Флойда-Уоршалла. Она печатает исходную весовую матрицу, итоговую матрицу кратчайших расстояний, матрицу восстановления пути и путь между выбранными пользователем вершинами. Если путь отсутствует или найден отрицательный цикл, функция выводит соответствующее предупреждение.

\textbf{Вход:} результат \texttt{res}; весовая матрица \texttt{W}; количество вершин \texttt{n}; начальная вершина \texttt{src}; конечная вершина \texttt{dst}.

\textbf{Выход:} подробный вывод таблиц и восстановленного кратчайшего пути в консоли.

\subsubsection*{Функции лабораторной работы №3}

\noindent\textbf{Функция} \coderef{std::vector<std::vector<int>> generateCapacityMatrix(const Graph& g)}{lst:lab3-ford-fulkerson-cpp}\\
Функция создаёт матрицу пропускных способностей для ориентированного графа. Для каждой существующей дуги генерируется целое значение пропускной способности, а для отсутствующих дуг остаётся ноль. Эта матрица используется в алгоритмах максимального потока и потока минимальной стоимости.

\textbf{Вход:} ориентированный граф \texttt{g}.

\textbf{Выход:} матрица пропускных способностей, где положительное число означает наличие дуги и её ёмкость.

\noindent\textbf{Функция} \coderef{std::vector<std::vector<int>> generateCostMatrix(const Graph& g)}{lst:lab3-ford-fulkerson-cpp}\\
Функция создаёт матрицу стоимостей для тех же дуг, которые присутствуют в графе. Стоимость показывает цену прохождения одной единицы потока по конкретной дуге.

\textbf{Вход:} ориентированный граф \texttt{g}.

\textbf{Выход:} матрица стоимостей, где ноль для отсутствующей дуги означает, что поток по ней не допускается.

\noindent\textbf{Функция} \coderef{FFResult fordFulkerson(int n, const std::vector<std::vector<int>>& cap, int src, int dst)}{lst:lab3-ford-fulkerson-cpp}\\
Функция реализует алгоритм Форда-Фалкерсона в варианте Эдмондса-Карпа. Это означает, что каждый увеличивающий путь в остаточной сети ищется с помощью поиска в ширину. После нахождения пути определяется минимальная остаточная пропускная способность на этом пути, затем поток увеличивается, а обратные рёбра в остаточной сети обновляются.

\textbf{Вход:} количество вершин \texttt{n}; матрица пропускных способностей \texttt{cap}; исток \texttt{src}; сток \texttt{dst}.

\textbf{Выход:} структура \texttt{FFResult}, содержащая итоговую матрицу потока, значение максимального потока, исток и сток.

\noindent\textbf{Функция} \coderef{FFResult fordFulkersonSilent(int n, const std::vector<std::vector<int>>& cap, int src, int dst)}{lst:lab3-ford-fulkerson-cpp}\\
Функция выполняет тот же расчёт максимального потока, что и \coderef{fordFulkerson}{lst:lab3-ford-fulkerson-cpp}, но не печатает промежуточные увеличивающие пути. Она используется внутри функции минимального стоимостного потока, где нужен только численный результат и матрица потока.

\textbf{Вход:} количество вершин \texttt{n}; матрица пропускных способностей \texttt{cap}; исток \texttt{src}; сток \texttt{dst}.

\textbf{Выход:} структура \texttt{FFResult} без промежуточного вывода в консоль.

\noindent\textbf{Функция} \coderef{void printCapacityMatrix(const std::vector<std::vector<int>>& cap, int n)}{lst:lab3-ford-fulkerson-cpp}\\
Функция печатает матрицу пропускных способностей. Она нужна для того, чтобы пользователь мог видеть исходные данные перед запуском алгоритма максимального потока.

\textbf{Вход:} матрица \texttt{cap}; количество вершин \texttt{n}.

\textbf{Выход:} таблица пропускных способностей в консоли.

\noindent\textbf{Функция} \coderef{void printCostMatrix(const std::vector<std::vector<int>>& cost, int n)}{lst:lab3-ford-fulkerson-cpp}\\
Функция печатает матрицу стоимостей дуг. Она используется после генерации матриц для лабораторной работы №3.

\textbf{Вход:} матрица стоимостей \texttt{cost}; количество вершин \texttt{n}.

\textbf{Выход:} таблица стоимостей в консоли.

\noindent\textbf{Функция} \coderef{void printFFResult(const FFResult& res, const std::vector<std::vector<int>>& cap, int n)}{lst:lab3-ford-fulkerson-cpp}\\
Функция выводит результат алгоритма Форда-Фалкерсона. Она показывает максимальный поток, дуги с положительным потоком и соотношение потока к пропускной способности.

\textbf{Вход:} результат \texttt{res}; матрица пропускных способностей \texttt{cap}; количество вершин \texttt{n}.

\textbf{Выход:} форматированный отчёт о максимальном потоке в консоли.

\noindent\textbf{Функция} \coderef{MCFResult minCostFlow(int n, const std::vector<std::vector<int>>& cap, const std::vector<std::vector<int>>& cost, int src, int dst)}{lst:lab3-min-cost-flow-cpp}\\
Функция вычисляет поток минимальной стоимости. Сначала с помощью \coderef{fordFulkersonSilent}{lst:lab3-ford-fulkerson-cpp} находится максимальный поток, затем целевой поток принимается равным $\lfloor 2/3 \cdot \varphi_{max} \rfloor$. После этого поток набирается последовательным добавлением кратчайших по стоимости путей в остаточной сети. Для каждого выбранного пути учитываются остаточные ёмкости и обратные дуги.

\textbf{Вход:} количество вершин \texttt{n}; матрица пропускных способностей \texttt{cap}; матрица стоимостей \texttt{cost}; исток \texttt{src}; сток \texttt{dst}.

\textbf{Выход:} структура \texttt{MCFResult} с максимальным потоком, целевым потоком, достигнутым потоком, итоговой стоимостью и матрицей потока.

\noindent\textbf{Функция} \coderef{void printMCFResult(const MCFResult& res, const std::vector<std::vector<int>>& cap, const std::vector<std::vector<int>>& cost, int n)}{lst:lab3-min-cost-flow-cpp}\\
Функция выводит результат минимального стоимостного потока. Она печатает значение максимального потока, требуемый поток, фактически найденный поток, список задействованных дуг и проверочный расчёт общей стоимости.

\textbf{Вход:} результат \texttt{res}; матрица пропускных способностей \texttt{cap}; матрица стоимостей \texttt{cost}; количество вершин \texttt{n}.

\textbf{Выход:} подробный вывод потока минимальной стоимости в консоли.

\subsubsection*{Функции лабораторной работы №4}

\noindent\textbf{Функция} \coderef{long long countSpanningTrees(const Graph& g)}{lst:lab4-kirchhoff-cpp}\\
Функция считает количество остовных деревьев с помощью матричной теоремы Кирхгофа. Сначала строится матрица Лапласа: на диагонали записываются степени вершин, а на местах существующих рёбер вне диагонали ставится \texttt{-1}. Затем удаляются одна строка и один столбец, после чего вычисляется определитель полученной матрицы.

\textbf{Вход:} неориентированный граф \texttt{g}.

\textbf{Выход:} количество остовных деревьев типа \texttt{long long}.

\noindent\textbf{Функция} \coderef{void printKirchhoffResult(const Graph& g, long long count)}{lst:lab4-kirchhoff-cpp}\\
Функция выводит результат применения теоремы Кирхгофа. Она показывает количество остовных деревьев для текущего графа.

\textbf{Вход:} граф \texttt{g}; найденное количество остовных деревьев \texttt{count}.

\textbf{Выход:} текстовый результат в консоли.

\noindent\textbf{Функция} \coderef{BoruvkaResult boruvka(const Graph& g, const std::vector<std::vector<double>>& W)}{lst:lab4-boruvka-cpp}\\
Функция реализует алгоритм Борувки для построения минимального остовного дерева. На каждой итерации для каждой компоненты связности выбирается самое дешёвое исходящее ребро. Эти рёбра добавляются в остов, а соответствующие компоненты объединяются. Процесс продолжается, пока не останется одна компонента или пока дальнейшее объединение невозможно.

\textbf{Вход:} неориентированный граф \texttt{g}; весовая матрица \texttt{W}.

\textbf{Выход:} структура \texttt{BoruvkaResult}, содержащая список рёбер минимального остова и его суммарный вес.

\noindent\textbf{Функция} \coderef{void printBoruvkaResult(const BoruvkaResult& res)}{lst:lab4-boruvka-cpp}\\
Функция печатает рёбра минимального остовного дерева и общую стоимость остова. Она используется после запуска алгоритма Борувки и перед алгоритмами, которым нужен построенный остов.

\textbf{Вход:} структура \texttt{BoruvkaResult}.

\textbf{Выход:} список рёбер остова и его вес в консоли.

\noindent\textbf{Функция} \coderef{EdgeCoverResult minEdgeCover(const Graph& g, const std::vector<std::vector<double>>& W, bool useMST, const std::vector<MSTEdge>& mstEdges)}{lst:lab4-edge-cover-cpp}\\
Функция строит минимальное рёберное покрытие. Если выбран режим \texttt{useMST}, покрытие строится по рёбрам минимального остова; иначе используется весь исходный граф. Сначала ищется максимальное паросочетание, затем для каждой непокрытой вершины добавляется любое подходящее инцидентное ребро. Для общего неориентированного графа используется алгоритм Эдмондса с сжатием цветков.

\textbf{Вход:} граф \texttt{g}; весовая матрица \texttt{W}; флаг \texttt{useMST}; список рёбер остова \texttt{mstEdges}.

\textbf{Выход:} структура \texttt{EdgeCoverResult}, содержащая рёбра покрытия и размер найденного паросочетания.

\noindent\textbf{Функция} \coderef{void printEdgeCoverResult(const EdgeCoverResult& res)}{lst:lab4-edge-cover-cpp}\\
Функция выводит результат построения рёберного покрытия. Для каждого ребра показываются его концы и вес, а также печатается размер паросочетания, использованного при построении покрытия.

\textbf{Вход:} структура \texttt{EdgeCoverResult}.

\textbf{Выход:} список рёбер минимального покрытия в консоли.

\noindent\textbf{Функция} \coderef{PruferResult pruferEncode(const std::vector<MSTEdge>& mstEdges, int n)}{lst:lab4-prufer-cpp}\\
Функция кодирует минимальное остовное дерево кодом Прюфера. На каждом шаге выбирается лист с минимальным номером, в код записывается его единственный сосед, затем лист удаляется из текущего дерева. Вес удалённого ребра сохраняется отдельно, чтобы при декодировании восстановить не только структуру дерева, но и веса рёбер.

\textbf{Вход:} список рёбер остова \texttt{mstEdges}; количество вершин \texttt{n}.

\textbf{Выход:} структура \texttt{PruferResult}, содержащая код Прюфера и список весов удаляемых рёбер.

\noindent\textbf{Функция} \coderef{std::vector<MSTEdge> pruferDecode(const std::vector<int>& code, const std::vector<double>& weights, int n)}{lst:lab4-prufer-cpp}\\
Функция восстанавливает дерево по коду Прюфера. На каждом шаге выбирается минимальный лист, которого нет в текущем коде, и соединяется с первым элементом кода. Вес ребра берётся из сохранённого массива \texttt{weights}. После обработки всех элементов кода соединяются две оставшиеся вершины.

\textbf{Вход:} код Прюфера \texttt{code}; список весов \texttt{weights}; количество вершин \texttt{n}.

\textbf{Выход:} список рёбер восстановленного дерева.

\noindent\textbf{Функция} \coderef{void printPruferResult(const PruferResult& res, int n)}{lst:lab4-prufer-cpp}\\
Функция выводит код Прюфера и сохранённые веса рёбер. Она показывает, какая последовательность получилась после кодирования остова.

\textbf{Вход:} результат кодирования \texttt{res}; количество вершин \texttt{n}.

\textbf{Выход:} код Прюфера и веса в консоли.

\noindent\textbf{Функция} \coderef{void printDecodedTree(const std::vector<MSTEdge>& edges)}{lst:lab4-prufer-cpp}\\
Функция выводит дерево, восстановленное из кода Прюфера. Это позволяет сравнить исходный остов и результат декодирования.

\textbf{Вход:} список восстановленных рёбер \texttt{edges}.

\textbf{Выход:} список рёбер декодированного дерева в консоли.

\subsubsection*{Функции лабораторной работы №5}

\noindent\textbf{Функция} \coderef{EulerianResult buildEulerianCycle(const Graph& g)}{lst:lab5-eulerian-cpp}\\
Функция проверяет, является ли неориентированный граф эйлеровым, и строит эйлеров цикл. Сначала вычисляются степени всех вершин и проверяется связность. Если граф уже связный и все степени чётные, сразу запускается построение цикла. Если условия не выполнены, программа пытается локально изменить копию графа: добавить или удалить рёбра так, чтобы сделать степени чётными и не нарушить связность. Исходный объект \texttt{Graph} при этом не изменяется.

\textbf{Вход:} неориентированный граф \texttt{g}.

\textbf{Выход:} структура \texttt{EulerianResult}, содержащая исходные степени, нечётные вершины, выполненные изменения и найденный эйлеров цикл.

\noindent\textbf{Функция} \coderef{void printEulerianResult(const EulerianResult& res)}{lst:lab5-eulerian-cpp}\\
Функция выводит результат проверки графа на эйлеровость. Она печатает исходные степени вершин, список нечётных вершин, информацию о добавленных и удалённых рёбрах, а также сам эйлеров цикл, если его удалось построить.

\textbf{Вход:} структура \texttt{EulerianResult}.

\textbf{Выход:} подробный отчёт об эйлеровом цикле в консоли.

\noindent\textbf{Функция} \coderef{std::vector<FundamentalCutset> buildFundamentalCutsets(const Graph& g, const std::vector<MSTEdge>& mstEdges)}{lst:lab5-fundamental-cutsets-cpp}\\
Функция строит фундаментальную систему разрезов относительно минимального остова. Для каждого ребра остова оно временно удаляется из дерева. После удаления дерево распадается на две компоненты, и функция ищет все рёбра исходного графа, которые соединяют эти две компоненты. Полученный набор рёбер образует фундаментальный разрез.

\textbf{Вход:} исходный граф \texttt{g}; список рёбер минимального остова \texttt{mstEdges}.

\textbf{Выход:} вектор структур \texttt{FundamentalCutset}, каждая из которых содержит ребро остова и соответствующий разрез.

\noindent\textbf{Функция} \coderef{void printFundamentalCutsets(const std::vector<FundamentalCutset>& cutsets)}{lst:lab5-fundamental-cutsets-cpp}\\
Функция выводит все построенные фундаментальные разрезы. Для каждого разреза показывается номер, соответствующее ребро остова и список рёбер, входящих в разрез.

\textbf{Вход:} список фундаментальных разрезов \texttt{cutsets}.

\textbf{Выход:} таблица фундаментальных разрезов в консоли.

\noindent\textbf{Функция} \coderef{void printCutsetSymmetricDifference(const std::vector<FundamentalCutset>& cutsets, const std::vector<int>& selectedIndices)}{lst:lab5-fundamental-cutsets-cpp}\\
Функция вычисляет симметрическую разность выбранных фундаментальных разрезов. Если ребро встречается в выбранных разрезах нечётное число раз, оно остаётся в результате; если чётное число раз, оно исключается. Это соответствует сложению множеств рёбер по модулю два.

\textbf{Вход:} список разрезов \texttt{cutsets}; номера выбранных разрезов \texttt{selectedIndices}.

\textbf{Выход:} множество рёбер, полученное в результате симметрической разности, выведенное в консоль.

\subsubsection*{Функции объединяющего меню}

\noindent\textbf{Функция} \coderef{static inline int D(int v)}{lst:lab-combined-main-cpp}\\
Функция переводит внутреннюю нумерацию вершин в пользовательскую. В программе вершины хранятся с нуля, но пользователю они показываются с единицы.

\textbf{Вход:} номер вершины \texttt{v} во внутренней нумерации.

\textbf{Выход:} номер вершины \texttt{v + 1} для вывода в консоль.

\noindent\textbf{Функция} \coderef{static void clearInput()}{lst:lab-combined-main-cpp}\\
Функция очищает поток ввода после неправильного или завершённого чтения. Она сбрасывает флаг ошибки и удаляет оставшиеся символы до конца строки.

\textbf{Вход:} стандартный поток ввода \texttt{std::cin}.

\textbf{Выход:} очищенное состояние потока ввода.

\noindent\textbf{Функция} \coderef{static int readInt(const std::string& prompt)}{lst:lab-combined-main-cpp}\\
Функция безопасно считывает целое число. Она выводит приглашение, пытается прочитать значение и повторяет запрос, если пользователь ввёл некорректные данные. Это используется во всех пунктах меню, где требуется ввод чисел.

\textbf{Вход:} строка-приглашение \texttt{prompt}.

\textbf{Выход:} корректно считанное целое число.

\noindent\textbf{Функции} \coderef{requireGraph()}{lst:lab-combined-main-cpp}, \coderef{requireWeight()}{lst:lab-combined-main-cpp}, \coderef{requireFlowMatrices()}{lst:lab-combined-main-cpp}, \coderef{requireMST()}{lst:lab-combined-main-cpp}\\
Эти функции выводят сообщения об ошибках, если пользователь пытается запустить алгоритм без необходимых предварительных данных. Например, весовую матрицу нельзя использовать до генерации графа, а код Прюфера нельзя построить до нахождения минимального остова.

\textbf{Вход:} текущее состояние глобальных переменных программы.

\textbf{Выход:} диагностическое сообщение в консоли.

\noindent\textbf{Функция} \coderef{static bool requireLab4Undirected()}{lst:lab-combined-main-cpp}\\
Функция проверяет, можно ли запускать алгоритмы лабораторной работы №4. Для этих алгоритмов требуется уже сгенерированный неориентированный граф.

\textbf{Вход:} состояние текущего графа.

\textbf{Выход:} \texttt{true}, если граф существует и является неориентированным; иначе \texttt{false} и сообщение об ошибке.

\noindent\textbf{Функция} \coderef{static bool requireLab5Undirected()}{lst:lab-combined-main-cpp}\\
Функция выполняет аналогичную проверку для лабораторной работы №5. Она не даёт запускать эйлеров цикл и фундаментальные разрезы на ориентированном графе.

\textbf{Вход:} состояние текущего графа.

\textbf{Выход:} \texttt{true} при выполнении условий запуска; иначе \texttt{false}.

\noindent\textbf{Функция} \coderef{void menuGenerateGraph()}{lst:lab-combined-main-cpp}\\
Функция реализует пункт меню генерации графа. Она считывает количество вершин и тип графа, вызывает \coderef{generateDAG}{lst:shared-dag-generator-cpp}, сохраняет результат в глобальную переменную и сбрасывает ранее построенные матрицы и результаты, потому что они уже не соответствуют новому графу.

\textbf{Вход:} количество вершин и тип графа, введённые пользователем.

\textbf{Выход:} новый граф в глобальной переменной \texttt{gGraph} и вывод матрицы смежности.

\noindent\textbf{Функция} \coderef{void menuGenerateWeightMatrix()}{lst:lab-combined-main-cpp}\\
Функция реализует пункт меню генерации весовой матрицы. Она проверяет наличие графа, считывает выбранный режим весов, вызывает \coderef{generateWeightMatrix}{lst:shared-weight-matrix-cpp} и выводит полученную матрицу.

\textbf{Вход:} текущий граф и выбранный режим весов.

\textbf{Выход:} глобальная матрица \texttt{gWeightMatrix} и её вывод в консоли.

\noindent\textbf{Функция} \coderef{static void printShimbellMatrix(const Matrix& M)}{lst:lab-combined-main-cpp}\\
Функция используется в меню метода Шимбелла для аккуратного вывода матриц минимальных и максимальных путей. Она отдельно обрабатывает значения бесконечности.

\textbf{Вход:} матрица \texttt{M}.

\textbf{Выход:} форматированный вывод матрицы.

\noindent\textbf{Функции} \coderef{menuEccentricity()}{lst:lab-combined-main-cpp}, \coderef{menuShimbell()}{lst:lab-combined-main-cpp}, \coderef{menuPaths()}{lst:lab-combined-main-cpp}\\
Эти функции связывают пункты меню лабораторной работы №1 с алгоритмами. Они проверяют наличие нужных данных, считывают параметры пользователя и вызывают соответственно \coderef{computeMetrics}{lst:lab1-eccentricity-cpp}, \coderef{shimbell}{lst:lab1-shimbell-cpp} и \coderef{findAllPaths}{lst:lab1-path-counter-cpp}.

\textbf{Вход:} текущий граф, весовая матрица при необходимости и параметры, введённые пользователем.

\textbf{Выход:} результаты алгоритмов лабораторной работы №1 в консоли.

\noindent\textbf{Функции} \coderef{menuDFS()}{lst:lab-combined-main-cpp} и \coderef{menuFloydWarshall()}{lst:lab-combined-main-cpp}\\
Эти функции отвечают за запуск алгоритмов лабораторной работы №2. Первая считывает стартовую вершину и запускает \coderef{dfs}{lst:lab2-dfs-cpp}. Вторая подготавливает весовую матрицу для Флойда-Уоршалла, устанавливая нули на диагонали, затем считывает начальную и конечную вершины и вызывает \coderef{floydWarshall}{lst:lab2-floyd-warshall-cpp}.

\textbf{Вход:} текущий граф, весовая матрица и выбранные вершины.

\textbf{Выход:} обход в глубину или кратчайший путь между выбранными вершинами.

\noindent\textbf{Функции} \coderef{menuGenerateFlowMatrices()}{lst:lab-combined-main-cpp}, \coderef{menuFordFulkerson()}{lst:lab-combined-main-cpp}, \coderef{menuMinCostFlow()}{lst:lab-combined-main-cpp}\\
Эти функции реализуют пункты меню лабораторной работы №3. Первая создаёт матрицы пропускных способностей и стоимостей, вторая запускает максимальный поток, третья запускает поток минимальной стоимости. Перед запуском выполняется проверка, что граф ориентированный и необходимые матрицы уже созданы.

\textbf{Вход:} ориентированный граф, выбранные исток и сток, матрицы пропускных способностей и стоимостей.

\textbf{Выход:} максимальный поток или поток минимальной стоимости, выведенный в консоль.

\noindent\textbf{Функции} \coderef{menuKirchhoff()}{lst:lab-combined-main-cpp}, \coderef{menuBoruvka()}{lst:lab-combined-main-cpp}, \coderef{menuEdgeCover()}{lst:lab-combined-main-cpp}, \coderef{menuPrufer()}{lst:lab-combined-main-cpp}\\
Эти функции запускают алгоритмы лабораторной работы №4. Они работают только с неориентированным графом. \coderef{menuKirchhoff}{lst:lab-combined-main-cpp} считает количество остовных деревьев, \coderef{menuBoruvka}{lst:lab-combined-main-cpp} строит минимальный остов, \coderef{menuEdgeCover}{lst:lab-combined-main-cpp} строит рёберное покрытие, а \coderef{menuPrufer}{lst:lab-combined-main-cpp} кодирует и декодирует остов кодом Прюфера.

\textbf{Вход:} неориентированный граф, весовая матрица, а также построенный остов для задач, которые от него зависят.

\textbf{Выход:} результаты алгоритмов лабораторной работы №4.

\noindent\textbf{Функции} \coderef{menuEulerianCycle()}{lst:lab-combined-main-cpp} и \coderef{menuFundamentalCutsets()}{lst:lab-combined-main-cpp}\\
Эти функции реализуют пункты меню лабораторной работы №5. Первая строит эйлеров цикл или показывает, почему это невозможно. Вторая строит минимальный остов, затем по нему формирует фундаментальные разрезы и позволяет выбрать разрезы для вычисления симметрической разности.

\textbf{Вход:} неориентированный граф, весовая матрица и выбранные пользователем номера разрезов.

\textbf{Выход:} эйлеров цикл либо система фундаментальных разрезов и их симметрическая разность.

\noindent\textbf{Функция} \coderef{void printMenu()}{lst:lab-combined-main-cpp}\\
Функция печатает главное меню программы. В меню перечислены все пункты: генерация графа, построение матриц, запуск алгоритмов лабораторных работ №1--№5 и выход.

\textbf{Вход:} входных параметров нет.

\textbf{Выход:} текст главного меню в консоли.

\noindent\textbf{Функция} \coderef{int main()}{lst:lab-combined-main-cpp}\\
Главная функция программы организует общий цикл работы приложения. Пока пользователь не выбрал пункт \texttt{0}, программа печатает меню, считывает номер команды через \coderef{readInt}{lst:lab-combined-main-cpp} и с помощью оператора \texttt{switch} вызывает соответствующую функцию меню. Такой подход позволяет объединить пять лабораторных работ в одном исполняемом файле.

\textbf{Вход:} команды пользователя из консоли.

\textbf{Выход:} выполнение выбранных алгоритмов и завершение программы после выбора пункта выхода.



\clearpage
\section{Результаты работы программы}

\subsection{Задания 1--11: ориентированный граф на 5 вершинах}

После запуска программы выводится главное меню, в котором доступны пункты всех пяти лабораторных работ. Для заданий 1--11 был сгенерирован ориентированный граф на 5 вершинах, так как первые три лабораторные работы используют ориентированный ациклический граф и алгоритмы для маршрутов, кратчайших путей и потоков. Главное меню программы показано на рис.~\ref{fig:res-ui}.

\begin{figure}[H]
	\centering
	\includegraphics[width=0.95\textwidth,height=0.78\textheight,keepaspectratio]{\detokenize{1 to 11 with directed graph vertices 5/ui.png}}
	\caption{Главное меню объединённой программы}
	\label{fig:res-ui}
\end{figure}

При выборе пункта <<1>> пользователь вводит количество вершин и тип графа. Для первой части результатов выбран ориентированный граф; программа строит матрицу смежности и сообщает, что для дальнейших алгоритмов необходимо сгенерировать весовую матрицу~--- рис.~\ref{fig:res-generate-directed}.

\begin{figure}[H]
	\centering
	\includegraphics[width=0.95\textwidth,height=0.78\textheight,keepaspectratio]{\detokenize{1 to 11 with directed graph vertices 5/1.generate graph directed Vertices 5.png}}
	\caption{Генерация ориентированного графа на 5 вершинах}
	\label{fig:res-generate-directed}
\end{figure}

При выборе пункта <<2>> выводится полное представление графа: список смежности, список дуг и матрица смежности. Это позволяет проверить структуру сгенерированного ориентированного графа перед запуском алгоритмов~--- рис.~\ref{fig:res-show-directed}.

\begin{figure}[H]
	\centering
	\includegraphics[width=0.95\textwidth,height=0.78\textheight,keepaspectratio]{\detokenize{1 to 11 with directed graph vertices 5/2. show graph V5.png}}
	\caption{Вывод списка смежности, списка дуг и матрицы смежности}
	\label{fig:res-show-directed}
\end{figure}

При выборе пункта <<3>> пользователь выбирает режим генерации весов: положительные, отрицательные или смешанные. После этого программа выводит весовую матрицу, где значение INF означает отсутствие дуги~--- рис.~\ref{fig:res-weight-directed}.

\begin{figure}[H]
	\centering
	\includegraphics[width=0.95\textwidth,height=0.78\textheight,keepaspectratio]{\detokenize{1 to 11 with directed graph vertices 5/3. generate weight matrix.png}}
	\caption{Генерация весовой матрицы для ориентированного графа}
	\label{fig:res-weight-directed}
\end{figure}

При выборе пункта <<4>> вычисляются расстояния между вершинами, эксцентриситеты, радиус, диаметр, центр графа и диаметральные вершины. Результат работы алгоритма представлен в виде матрицы расстояний и итоговой таблицы~--- рис.~\ref{fig:res-eccentricity}.

\begin{figure}[H]
	\centering
	\includegraphics[width=0.95\textwidth,height=0.78\textheight,keepaspectratio]{\detokenize{1 to 11 with directed graph vertices 5/4. Eccentricities, center, diametral vertices.png}}
	\caption{Эксцентриситеты, центр и диаметральные вершины графа}
	\label{fig:res-eccentricity}
\end{figure}

При выборе пункта <<5>> пользователь задаёт количество шагов $k$, после чего программа применяет метод Шимбелла. В результате выводятся матрицы минимальных и максимальных путей, состоящих ровно из заданного количества дуг~--- рис.~\ref{fig:res-shimbell}.

\begin{figure}[H]
	\centering
	\includegraphics[width=0.95\textwidth,height=0.78\textheight,keepaspectratio]{\detokenize{1 to 11 with directed graph vertices 5/5. shimbell's method.png}}
	\caption{Метод Шимбелла для ориентированного взвешенного графа}
	\label{fig:res-shimbell}
\end{figure}

При выборе пункта <<6>> пользователь вводит начальную и конечную вершины. Программа определяет, существует ли маршрут между ними, подсчитывает количество простых маршрутов и выводит найденные маршруты~--- рис.~\ref{fig:res-routes}.

\begin{figure}[H]
	\centering
	\includegraphics[width=0.95\textwidth,height=0.78\textheight,keepaspectratio]{\detokenize{1 to 11 with directed graph vertices 5/6. find all routes between two vertices.png}}
	\caption{Поиск всех маршрутов между двумя вершинами}
	\label{fig:res-routes}
\end{figure}

При выборе пункта <<7>> выполняется обход графа поиском в глубину из заданной вершины. В результате выводится порядок обхода, количество посещённых вершин и число итераций алгоритма~--- рис.~\ref{fig:res-dfs}.

\begin{figure}[H]
	\centering
	\includegraphics[width=0.95\textwidth,height=0.72\textheight,keepaspectratio]{\detokenize{1 to 11 with directed graph vertices 5/7. DFS traversal.png}}
	\caption{Обход графа поиском в глубину}
	\label{fig:res-dfs}
\end{figure}

При выборе пункта <<8>> пользователь вводит начальную и конечную вершины для поиска кратчайшего пути. Алгоритм Флойда-Уоршалла выводит исходную весовую матрицу, матрицу кратчайших расстояний, матрицу восстановления пути, сам путь и количество итераций~--- рис.~\ref{fig:res-floyd}.

\begin{figure}[H]
	\centering
	\includegraphics[width=0.95\textwidth,height=0.78\textheight,keepaspectratio]{\detokenize{1 to 11 with directed graph vertices 5/8. Shortest path Floyd-warshall.png}}
	\caption{Поиск кратчайшего пути алгоритмом Флойда-Уоршалла}
	\label{fig:res-floyd}
\end{figure}

При выборе пункта <<9>> на основе ориентированного графа генерируются матрицы пропускных способностей и стоимостей. Дополнительно программа выводит входящие и исходящие степени вершин, чтобы пользователю было проще выбрать исток и сток~--- рис.~\ref{fig:res-flow-matrices}.

\begin{figure}[H]
	\centering
	\includegraphics[width=0.95\textwidth,height=0.78\textheight,keepaspectratio]{\detokenize{1 to 11 with directed graph vertices 5/9. generate capacity and cost matrices.png}}
	\caption{Генерация матриц пропускных способностей и стоимостей}
	\label{fig:res-flow-matrices}
\end{figure}

При выборе пункта <<10>> программа находит максимальный поток алгоритмом Форда-Фалкерсона. На экран выводятся увеличивающие пути, значение максимального потока и распределение потока по дугам~--- рис.~\ref{fig:res-ford-fulkerson}.

\begin{figure}[H]
	\centering
	\includegraphics[width=0.95\textwidth,height=0.78\textheight,keepaspectratio]{\detokenize{1 to 11 with directed graph vertices 5/10. Max flow ford fulkerson.png}}
	\caption{Нахождение максимального потока алгоритмом Форда-Фалкерсона}
	\label{fig:res-ford-fulkerson}
\end{figure}

При выборе пункта <<11>> вычисляется поток минимальной стоимости. Величина требуемого потока берётся равной $\lfloor 2/3 \cdot \varphi_{\max} \rfloor$. Программа выводит найденные кратчайшие по стоимости пути, достигнутый поток и итоговую стоимость~--- рис.~\ref{fig:res-min-cost-flow}.

\begin{figure}[H]
	\centering
	\includegraphics[width=0.95\textwidth,height=0.78\textheight,keepaspectratio]{\detokenize{1 to 11 with directed graph vertices 5/11. Min-cost flow.png}}
	\caption{Вычисление потока минимальной стоимости}
	\label{fig:res-min-cost-flow}
\end{figure}

\clearpage
\subsection{Задания 12--17: неориентированный граф на 7 вершинах}

Для заданий 12--17 был сгенерирован неориентированный граф на 7 вершинах, так как алгоритмы четвёртой и пятой лабораторных работ работают с неориентированным графом. После генерации выводится матрица смежности~--- рис.~\ref{fig:res-generate-undirected}.

\begin{figure}[H]
	\centering
	\includegraphics[width=0.95\textwidth,height=0.78\textheight,keepaspectratio]{\detokenize{12 to 17 with undirected graph vertices 7/1. generate graph for 12 to 17 lab 4 and 5.png}}
	\caption{Генерация неориентированного графа на 7 вершинах}
	\label{fig:res-generate-undirected}
\end{figure}

При выборе пункта <<12>> программа строит матрицу Кирхгофа, удаляет первую строку и первый столбец для получения алгебраического дополнения и вычисляет определитель. Полученное значение является числом остовных деревьев графа~--- рис.~\ref{fig:res-kirchhoff}.

\begin{figure}[H]
	\centering
	\includegraphics[width=0.95\textwidth,height=0.78\textheight,keepaspectratio]{\detokenize{12 to 17 with undirected graph vertices 7/12. kirchhoff count spanning trees.png}}
	\caption{Подсчёт числа остовных деревьев по теореме Кирхгофа}
	\label{fig:res-kirchhoff}
\end{figure}

При выборе пункта <<13>> строится минимальное остовное дерево алгоритмом Борувки. Программа по шагам выводит компоненты связности, добавляемые безопасные рёбра и итоговый вес полученного остова~--- рис.~\ref{fig:res-boruvka}.

\begin{figure}[H]
	\centering
	\includegraphics[width=0.95\textwidth,height=0.78\textheight,keepaspectratio]{\detokenize{12 to 17 with undirected graph vertices 7/13 Boruvka minimun spanning tree.png}}
	\caption{Построение минимального остова алгоритмом Борувки}
	\label{fig:res-boruvka}
\end{figure}

При выборе пункта <<14>> пользователь выбирает, искать минимальное рёберное покрытие на исходном графе или на построенном остове. В результате выводится мощность максимального паросочетания и рёбра минимального покрытия~--- рис.~\ref{fig:res-edge-cover}.

\begin{figure}[H]
	\centering
	\includegraphics[width=0.95\textwidth,height=0.78\textheight,keepaspectratio]{\detokenize{12 to 17 with undirected graph vertices 7/14. Minimum edge cover.png}}
	\caption{Нахождение минимального рёберного покрытия}
	\label{fig:res-edge-cover}
\end{figure}

При выборе пункта <<15>> построенный минимальный остов кодируется кодом Прюфера с сохранением весов, после чего выполняется обратное декодирование. На экран выводятся удаляемые листья, код Прюфера, веса рёбер и восстановленное дерево~--- рис.~\ref{fig:res-prufer}.

\begin{figure}[H]
	\centering
	\includegraphics[width=0.95\textwidth,height=0.78\textheight,keepaspectratio]{\detokenize{12 to 17 with undirected graph vertices 7/15. prufer encode  decode.png}}
	\caption{Кодирование и декодирование остова кодом Прюфера}
	\label{fig:res-prufer}
\end{figure}

При выборе пункта <<16>> программа проверяет, является ли граф эйлеровым. Если граф не удовлетворяет условию чётности степеней, выполняется модификация рёбер; после этого строится эйлеров цикл~--- рис.~\ref{fig:res-eulerian}.

\begin{figure}[H]
	\centering
	\includegraphics[width=0.95\textwidth,height=0.78\textheight,keepaspectratio]{\detokenize{12 to 17 with undirected graph vertices 7/16. Eulerian check modify .png}}
	\caption{Проверка графа на эйлеровость, модификация и построение эйлерова цикла}
	\label{fig:res-eulerian}
\end{figure}

При выборе пункта <<17>> программа строит минимальный остов алгоритмом Борувки, затем по каждому ребру остова получает фундаментальный разрез. После вывода фундаментальной системы пользователь выбирает номера разрезов, а программа вычисляет их симметрическую разность~--- рис.~\ref{fig:res-cutsets-1} и рис.~\ref{fig:res-cutsets-2}.

\begin{figure}[H]
	\centering
	\includegraphics[width=0.95\textwidth,height=0.78\textheight,keepaspectratio]{\detokenize{12 to 17 with undirected graph vertices 7/17.1 fundamental cutsets form MST.png}}
	\caption{Построение фундаментальной системы разрезов}
	\label{fig:res-cutsets-1}
\end{figure}

\begin{figure}[H]
	\centering
	\includegraphics[width=0.95\textwidth,height=0.78\textheight,keepaspectratio]{\detokenize{12 to 17 with undirected graph vertices 7/17.2 fundamental cutsets form MST.png}}
	\caption{Симметрическая разность выбранных фундаментальных разрезов}
	\label{fig:res-cutsets-2}
\end{figure}

\clearpage
\subsection{Обработка ошибочных ситуаций}

В программе предусмотрены проверки корректности пользовательских действий. Если пользователь пытается вызвать алгоритм без предварительно сгенерированного графа или без весовой матрицы, программа выводит сообщение с указанием нужного пункта меню~--- рис.~\ref{fig:res-error-1}.

\begin{figure}[H]
	\centering
	\includegraphics[width=0.9\textwidth,height=0.45\textheight,keepaspectratio]{\detokenize{error handling/error handle1.png}}
	\caption{Сообщение об ошибке при отсутствии необходимых данных}
	\label{fig:res-error-1}
\end{figure}

Для алгоритмов четвёртой и пятой лабораторных работ дополнительно проверяется тип графа. Если пользователь запускает алгоритм, требующий неориентированный граф, для ориентированного графа, программа сообщает о необходимости сгенерировать неориентированный граф~--- рис.~\ref{fig:res-error-2}.

\begin{figure}[H]
	\centering
	\includegraphics[width=0.9\textwidth,height=0.45\textheight,keepaspectratio]{\detokenize{error handling/error handle2 for 16 eulerian check.png}}
	\caption{Обработка неверного типа графа для алгоритма Эйлера}
	\label{fig:res-error-2}
\end{figure}

Также проверяется корректность номеров вершин и допустимость выбора истока и стока. При некорректном вводе программа не завершает работу аварийно, а выводит понятное сообщение и возвращает пользователя в меню~--- рис.~\ref{fig:res-error-3}.

\begin{figure}[H]
	\centering
	\includegraphics[width=0.9\textwidth,height=0.45\textheight,keepaspectratio]{\detokenize{error handling/error handle3 .png}}
	\caption{Обработка некорректного пользовательского ввода}
	\label{fig:res-error-3}
\end{figure}

\clearpage
\section*{Заключение}
\addcontentsline{toc}{section}{Заключение}

По результатам выполнения лабораторных работ можно сделать следующие выводы:

В первой лабораторной работе был реализован алгоритм генерации случайного связного ациклического графа с использованием нормального распределения. Для полученного графа были вычислены эксцентриситеты вершин, радиус, диаметр, центр и диаметральные вершины. Также был реализован метод Шимбелла для поиска минимальных и максимальных путей заданной длины и алгоритм определения всех маршрутов между двумя выбранными вершинами.

Во второй лабораторной работе был выполнен обход графа поиском в глубину. Для взвешенного графа был реализован алгоритм Флойда-Уоршалла, позволяющий находить кратчайшие пути между всеми парами вершин и восстанавливать путь между выбранными пользователем вершинами. В программе также выводится количество итераций, что позволяет оценивать трудоёмкость выполненных алгоритмов.

В третьей лабораторной работе на основе ориентированного графа были сгенерированы матрицы пропускных способностей и стоимостей. Для построенной транспортной сети был найден максимальный поток алгоритмом Форда-Фалкерсона. После этого был вычислен поток минимальной стоимости для величины, равной $\lfloor 2/3 \cdot \varphi_{\max} \rfloor$, с использованием ранее реализованного алгоритма Флойда-Уоршалла для поиска кратчайших по стоимости путей.

В четвёртой лабораторной работе для неориентированного графа было найдено число остовных деревьев с помощью матричной теоремы Кирхгофа. Минимальное остовное дерево было построено алгоритмом Борувки. Полученный остов был закодирован кодом Прюфера и затем декодирован с сохранением весов рёбер. Также был реализован поиск минимального рёберного покрытия как для исходного графа, так и для построенного остова.

В пятой лабораторной работе была выполнена проверка графа на эйлеровость. Если граф не являлся эйлеровым, программа выполняла его модификацию с выводом внесённых изменений и затем строила эйлеров цикл. Для чётного варианта была реализована фундаментальная система разрезов на основе минимального остова, а также получение других разрезов с помощью операции симметрической разности выбранных фундаментальных разрезов.

Выполнение лабораторных работ позволило изучить и реализовать основные алгоритмы теории графов для ориентированных и неориентированных графов: обходы, кратчайшие пути, потоки, остовы, покрытия, кодирование деревьев, эйлеровы циклы и фундаментальные разрезы. Все алгоритмы объединены в одну консольную программу с меню, что позволяет последовательно генерировать граф, строить необходимые матрицы и применять к ним алгоритмы из разных лабораторных работ.

Достоинством программы является модульная структура: общие сущности графа, генерации распределения и весовой матрицы вынесены в отдельные файлы, а алгоритмы разделены по лабораторным работам. Использование контейнеров STL, в частности \texttt{vector}, упрощает работу с матрицами и списками рёбер. В программе предусмотрены проверки пользовательского ввода и ограничения на запуск алгоритмов, например требование ориентированного графа для потоков и неориентированного графа для остовов, эйлерова цикла и разрезов.

К недостаткам можно отнести отсутствие графического представления графа: результаты выводятся в консоль в виде матриц, списков рёбер и текстовых протоколов работы алгоритмов. Кроме того, некоторые алгоритмы имеют не самую оптимальную асимптотику для больших графов, например поиск потока минимальной стоимости через повторное применение алгоритма Флойда-Уоршалла. В дальнейшем программу можно расширить визуализацией графов и добавить более эффективные реализации для работы с графами большого размера.

Лабораторные работы были выполнены на языке C++ в среде разработки CLion.

\clearpage
\addcontentsline{toc}{section}{Список использованной литературы}
\begin{thebibliography}{99}

	\bibitem{novikov}
	Новиков Ф.\,А. Дискретная математика для программистов~/ Ф.\,А.~Новиков.~---
	3-е изд.~--- СПб.:~Питер, 2009.~--- 364~с.

	\bibitem{shaporev}
	Шапорев С.\,Д. Дискретная математика~/ С.\,Д.~Шапорев.~---
	СПб.:~БХВ-Петербург, 2010.~--- 400~с.

	\bibitem{romanovsky}
	Романовский И.\,В. Дискретный анализ~/ И.\,В.~Романовский.~---
	3-е изд.~--- СПб.:~Невский диалект; БХВ-Петербург, 2003.~--- 320~с.

	\bibitem{haggarty}
	Хаггарти Р. Дискретная математика для программистов~/ Р.~Хаггарти; пер. с англ.~---
	М.:~Техносфера, 2005.~--- 326~с.

	\bibitem{cormen}
	Кормен Т., Лейзерсон Ч., Ривест Р., Штайн К. Алгоритмы. Построение и анализ~/ пер. с англ.~---
	3-е изд.~--- М.:~Вильямс, 2013.~--- 1328~с.

	\bibitem{asanov}
	Асанов М.\,О., Баранский В.\,А., Расин В.\,В. Дискретная математика: графы, матроиды, алгоритмы~/ М.\,О.~Асанов, В.\,А.~Баранский, В.\,В.~Расин.~---
	2-е изд.~--- СПб.:~Лань, 2010.~--- 368~с.

	\bibitem{vadzinsky}
	Вадзинский Р.\,Н. Справочник по вероятностным распределениям~/ Р.\,Н.~Вадзинский.~---
	СПб.:~Наука, 2001.~--- 295~с.

	\bibitem{erickson}
	Эрикссон Дж. Алгоритмы~/ Дж.~Эрикссон; пер. с англ.~---
	М.:~Вильямс, 2022.~--- 960~с.

	\bibitem{telematika}
	Секция «Телематика»~/ текст: электронный.~---
	URL: \url{https://tema.spbstu.ru/tgraph/}
	(Дата обращения: 20.05.2026)

\end{thebibliography}


\clearpage
\section*{Приложение А. Общие модули и объединяющая программа}
\addcontentsline{toc}{section}{Приложение А. Общие модули и объединяющая программа}

В приложении приведён полный исходный код общих модулей, используемых всеми лабораторными работами, а также код объединяющей консольной программы с меню.

\lstinputlisting[language=C++, caption={Файл constants.h}, label={lst:shared-constants-h}]{../graph-theory-labs/shared/constants.h}

\lstinputlisting[language=C++, caption={Файл graph.h}, label={lst:shared-graph-h}]{../graph-theory-labs/shared/graph.h}

\lstinputlisting[language=C++, caption={Файл graph.cpp}, label={lst:shared-graph-cpp}]{../graph-theory-labs/shared/graph.cpp}

\lstinputlisting[language=C++, caption={Файл distribution.h}, label={lst:shared-distribution-h}]{../graph-theory-labs/shared/distribution.h}

\lstinputlisting[language=C++, caption={Файл distribution.cpp}, label={lst:shared-distribution-cpp}]{../graph-theory-labs/shared/distribution.cpp}

\lstinputlisting[language=C++, caption={Файл dag\_generator.h}, label={lst:shared-dag-generator-h}]{../graph-theory-labs/shared/dag_generator.h}

\lstinputlisting[language=C++, caption={Файл dag\_generator.cpp}, label={lst:shared-dag-generator-cpp}]{../graph-theory-labs/shared/dag_generator.cpp}

\lstinputlisting[language=C++, caption={Файл weight\_matrix.h}, label={lst:shared-weight-matrix-h}]{../graph-theory-labs/shared/weight_matrix.h}

\lstinputlisting[language=C++, caption={Файл weight\_matrix.cpp}, label={lst:shared-weight-matrix-cpp}]{../graph-theory-labs/shared/weight_matrix.cpp}

\lstinputlisting[language=C++, caption={Файл main.cpp объединяющей программы}, label={lst:lab-combined-main-cpp}]{../graph-theory-labs/lab_combined/main.cpp}

\clearpage
\section*{Приложение Б. Исходный код лабораторной работы №1}
\addcontentsline{toc}{section}{Приложение Б. Исходный код лабораторной работы №1}

\lstinputlisting[language=C++, caption={Файл eccentricity.h}, label={lst:lab1-eccentricity-h}]{../graph-theory-labs/lab1/eccentricity.h}

\lstinputlisting[language=C++, caption={Файл eccentricity.cpp}, label={lst:lab1-eccentricity-cpp}]{../graph-theory-labs/lab1/eccentricity.cpp}

\lstinputlisting[language=C++, caption={Файл shimbell.h}, label={lst:lab1-shimbell-h}]{../graph-theory-labs/lab1/shimbell.h}

\lstinputlisting[language=C++, caption={Файл shimbell.cpp}, label={lst:lab1-shimbell-cpp}]{../graph-theory-labs/lab1/shimbell.cpp}

\lstinputlisting[language=C++, caption={Файл path\_counter.h}, label={lst:lab1-path-counter-h}]{../graph-theory-labs/lab1/path_counter.h}

\lstinputlisting[language=C++, caption={Файл path\_counter.cpp}, label={lst:lab1-path-counter-cpp}]{../graph-theory-labs/lab1/path_counter.cpp}

\clearpage
\section*{Приложение В. Исходный код лабораторной работы №2}
\addcontentsline{toc}{section}{Приложение В. Исходный код лабораторной работы №2}

\lstinputlisting[language=C++, caption={Файл dfs.h}, label={lst:lab2-dfs-h}]{../graph-theory-labs/lab2/dfs.h}

\lstinputlisting[language=C++, caption={Файл dfs.cpp}, label={lst:lab2-dfs-cpp}]{../graph-theory-labs/lab2/dfs.cpp}

\lstinputlisting[language=C++, caption={Файл floyd\_warshall.h}, label={lst:lab2-floyd-warshall-h}]{../graph-theory-labs/lab2/floyd_warshall.h}

\lstinputlisting[language=C++, caption={Файл floyd\_warshall.cpp}, label={lst:lab2-floyd-warshall-cpp}]{../graph-theory-labs/lab2/floyd_warshall.cpp}

\clearpage
\section*{Приложение Г. Исходный код лабораторной работы №3}
\addcontentsline{toc}{section}{Приложение Г. Исходный код лабораторной работы №3}

\lstinputlisting[language=C++, caption={Файл ford\_fulkerson.h}, label={lst:lab3-ford-fulkerson-h}]{../graph-theory-labs/lab3/ford_fulkerson.h}

\lstinputlisting[language=C++, caption={Файл ford\_fulkerson.cpp}, label={lst:lab3-ford-fulkerson-cpp}]{../graph-theory-labs/lab3/ford_fulkerson.cpp}

\lstinputlisting[language=C++, caption={Файл min\_cost\_flow.h}, label={lst:lab3-min-cost-flow-h}]{../graph-theory-labs/lab3/min_cost_flow.h}

\lstinputlisting[language=C++, caption={Файл min\_cost\_flow.cpp}, label={lst:lab3-min-cost-flow-cpp}]{../graph-theory-labs/lab3/min_cost_flow.cpp}

\clearpage
\section*{Приложение Д. Исходный код лабораторной работы №4}
\addcontentsline{toc}{section}{Приложение Д. Исходный код лабораторной работы №4}

\lstinputlisting[language=C++, caption={Файл kirchhoff.h}, label={lst:lab4-kirchhoff-h}]{../graph-theory-labs/lab4/kirchhoff.h}

\lstinputlisting[language=C++, caption={Файл kirchhoff.cpp}, label={lst:lab4-kirchhoff-cpp}]{../graph-theory-labs/lab4/kirchhoff.cpp}

\lstinputlisting[language=C++, caption={Файл boruvka.h}, label={lst:lab4-boruvka-h}]{../graph-theory-labs/lab4/boruvka.h}

\lstinputlisting[language=C++, caption={Файл boruvka.cpp}, label={lst:lab4-boruvka-cpp}]{../graph-theory-labs/lab4/boruvka.cpp}

\lstinputlisting[language=C++, caption={Файл edge\_cover.h}, label={lst:lab4-edge-cover-h}]{../graph-theory-labs/lab4/edge_cover.h}

\lstinputlisting[language=C++, caption={Файл edge\_cover.cpp}, label={lst:lab4-edge-cover-cpp}]{../graph-theory-labs/lab4/edge_cover.cpp}

\lstinputlisting[language=C++, caption={Файл prufer.h}, label={lst:lab4-prufer-h}]{../graph-theory-labs/lab4/prufer.h}

\lstinputlisting[language=C++, caption={Файл prufer.cpp}, label={lst:lab4-prufer-cpp}]{../graph-theory-labs/lab4/prufer.cpp}

\clearpage
\section*{Приложение Е. Исходный код лабораторной работы №5}
\addcontentsline{toc}{section}{Приложение Е. Исходный код лабораторной работы №5}

\lstinputlisting[language=C++, caption={Файл eulerian.h}, label={lst:lab5-eulerian-h}]{../graph-theory-labs/lab5/eulerian.h}

\lstinputlisting[language=C++, caption={Файл eulerian.cpp}, label={lst:lab5-eulerian-cpp}]{../graph-theory-labs/lab5/eulerian.cpp}

\lstinputlisting[language=C++, caption={Файл fundamental\_cutsets.h}, label={lst:lab5-fundamental-cutsets-h}]{../graph-theory-labs/lab5/fundamental_cutsets.h}

\lstinputlisting[language=C++, caption={Файл fundamental\_cutsets.cpp}, label={lst:lab5-fundamental-cutsets-cpp}]{../graph-theory-labs/lab5/fundamental_cutsets.cpp}


\end{document}
